

\chapter{The \texorpdfstring{$(\oo,1)$} --category of cobordisms}


The main goal of this thesis is to produce an $(\oo,1)$-category of stratified cobordisms with defects, and to sketch a higher version thereof. Successfully achieving this will require exploring a toolkit which can feel quite ornate and unfamiliar to a large swath of practicing mathematicians. To improve on this current state of affairs, introduce important ideas and results, and build a consistent narrative, this section is dedicated to introducing the two constructions we aim to generalize: the \emph{unstratified} $(\oo,n)$-category of cobordisms and the ordinary symmetric monoidal category of cobordisms with defects. 

\medskip

\section{A new construction of the \texorpdfstring{$(\oo,1)$}{(infty,1)}-bordism category}


In this section we present a new construction of the $d$-dimensional bordism
$(\oo,1)$-category. This construction is a spiritual cousin of earlier work:
most notably that of Galatius--Madsen--Tillman--Weiss \cite{GMTW09}, later
developed by Lurie \cite{Lur09}, and more recently by Scheimbauer--Calaque
\cite{CS19}. 

The pioneering work of Galatius--Madsen--Tillman--Weiss offered the first tools to construct an $(\oo,1)$-category of cobordisms. In their paper on the homotopy type of the cobordism category, they construct a topological $d$-dimensional bordism category by using classifying spaces. Following their work, Lurie, in his sketch of a proof of the cobordism hypothesis \cite{Lur09}, put forward a first construction of a model for an $(\oo,n)$-category of cobordisms using $n$-fold complete Segal spaces. This contained a small mistake preventing it from satisfying the \emph{essential constancy} condition. In order to address this, in her PhD thesis and in a subsequent paper with Damien Calaque, Scheimbauer--Calaque propose an alternative definition, which fixes Lurie's while retaining the use of $n$-fold complete Segal spaces. Despite this almost successful attempt at a \emph{corrigendum}, the work of Scheimbauer--Calaque contains a separate flaw when it claims to prove that the simplicial sets which model the levels of its $n$-fold Segal spaces are Kan complexes---which translates to a loss of fibrancy. We will explore this in the next chapter when we discuss the higher $n$ case.

Our approach will be to build $(\oo,n)$-categories using an alternative model, related to that of $n$-fold complete Segal spaces but more adapted to the construction of higher categories of cobordisms. For $n=1$, the case considered in this section, this means constructing an interval Segal space, whose theory we developed in \Cref{ch2:: title}. In the course of this construction we will pick up some tools from the theory of sheaves on manifolds, in particular as they are used and developed in the works of Berwick-Evans--Boavida de Brito--Pavlov \cite{BBEP24} and Ayala--Francis--Rozenblyum \cite{AFR1}. 

\subsection{A recollection of previous constructions}

Before getting into the deeper technical aspects of the mathematical work in this chapter, it is useful to first place it in its historical context. This serves both as a pedagogical aid for future users of higher cobordism categories and as a warm-up to the ideas of Calaque--Scheimbauer following Lurie, whose interpretation of cobordisms is, in spirit, different from the one found in earlier work on the subject. 

Instead of introducing the work of Lurie, as a means to point out what the work of Calaque--Scheimbauer fixes, we will instead present the main ideas driving the work and point out how it differs from Lurie and what fails in the original paper by the same author on the \emph{Cobordism Hypothesis}.

We are used to formulating cobordism categories from the \emph{bottom up}. In other words, we often define the cobordism category as a category whose $k$-morphisms are $(d-k)$-dimensional manifolds with corners and whose composition is given by gluing along the ``right'' boundary components. While this interpretation is retained in the works of Lurie and Calaque--Scheimbauer, their approach can be described as \emph{top down}. One of the most contrived aspects of defining higher categories of bordisms is composition, since it comes with a compounding pile of choices whose compatibility is easy to check, but very hard to write down.

This chapter treats only the $(\oo,1)$ case. We begin by discussing the work of
Lurie and Calaque--Scheimbauer, both to explain the background and to motivate
our choice of interval Segal spaces. 

For every $k\geq 0$, the space $(\bord_d^{\sf un})_k$ should encode the classifying space for composable chains of $k$ unoriented bordisms:
\begin{equation}
    \Sigma_0 \xra{M_1} \Sigma_1 \xra{M_2} \cdots \xra{M_k} \Sigma_k.
\end{equation}
Lurie's proposal is to define this space by picking auxiliary data and then showing it does not alter the relevant homotopy type. In this case, the auxiliary data is a finite-dimensional real vector space $V$ and a smooth embedding
\begin{equation}
    M \hookra V\times \RR
\end{equation}  
of the composed $d$-dimensional bordism $M = M_1 \sqcup_{\Sigma_1} \cdots \sqcup_{\Sigma_{k-1}} M_k$. Once we fix this embedding of $M$, we can encode the gluing regions i.e. the $\Sigma_i$'s, as certain intersections or ``cuts'' by transverse hyperplanes $V\times \{t_i\}$. To keep this visually grounded, let us consider the following example
\begin{equation}
    \tikzfig{figures/Ch3/BOR6}{\begin{tikzpicture}
	\begin{pgfonlayer}{nodelayer}
		\node [style=none] (0) at (-9.5, 0) {};
		\node [style=none] (1) at (-4.5, 3) {};
		\node [style=none] (2) at (-4.5, -3) {};
		\node [style=none] (3) at (9.5, 0) {};
		\node [style=none] (4) at (-4.5, 3) {};
		\node [style=none] (5) at (-4.5, -3) {};
		\node [style=none] (6) at (-7, 0.25) {};
		\node [style=none] (7) at (-2, 0.25) {};
		\node [style=none] (8) at (-6.75, -0.35) {};
		\node [style=none] (9) at (-2.5, -0.425) {};
		\node [style=none] (10) at (1.5, 0.25) {};
		\node [style=none] (11) at (6.5, 0.25) {};
		\node [style=none] (12) at (1.75, -0.35) {};
		\node [style=none] (13) at (6, -0.425) {};
		\node [style=none] (14) at (-12, 7) {};
		\node [style=none] (15) at (-12, -6.5) {};
		\node [style=none] (16) at (12, -6) {};
		\node [style=none] (17) at (-12.5, -6) {};
		\node [style=none] (18) at (-13, 7) {$V$};
		\node [style=none] (19) at (12, -6.75) {};
		\node [style=none] (20) at (-6.75, 6.5) {};
		\node [style=none] (21) at (-3.5, 9) {};
		\node [style=none] (22) at (-6.75, -7.5) {};
		\node [style=none] (23) at (-3.5, -5) {};
		\node [style=none] (32) at (-4.75, -6) {};
		\node [style=none] (33) at (-2, 6.5) {};
		\node [style=none] (34) at (1.25, 9) {};
		\node [style=none] (35) at (-2, -7.5) {};
		\node [style=none] (36) at (1.25, -5) {};
		\node [style=none] (37) at (2.25, 6.5) {};
		\node [style=none] (38) at (5.5, 9) {};
		\node [style=none] (39) at (2.25, -7.5) {};
		\node [style=none] (40) at (5.5, -5) {};
		\node [style=none] (41) at (0, -6) {};
		\node [style=none] (42) at (4.25, -6) {};
		\node [style=none] (43) at (-4.75, -7) {$t_0$};
		\node [style=none] (44) at (0, -7) {$t_1$};
		\node [style=none] (45) at (4.25, -7) {$t_2$};
		\node [style=none] (46) at (-4.75, -3) {};
		\node [style=none] (47) at (-4.75, -0.85) {};
		\node [style=none] (48) at (-4.75, 0.55) {};
		\node [style=none] (49) at (-4.75, 3) {};
		\node [style=none] (50) at (4.25, -3.65) {};
		\node [style=none] (51) at (4.25, -0.85) {};
		\node [style=none] (52) at (4.25, 0.55) {};
		\node [style=none] (53) at (4.325, 3.625) {};
		\node [style=none] (54) at (0, -3.325) {};
		\node [style=none] (55) at (-0.05, 3.325) {};
	\end{pgfonlayer}
	\begin{pgfonlayer}{edgelayer}
		\draw [style=geometric, in=90, out=-180] (1.center) to (0.center);
		\draw [style=geometric, in=180, out=-90] (0.center) to (2.center);
		\draw [style=geometric, in=90, out=0] (4.center) to (3.center);
		\draw [style=geometric, in=0, out=-90] (3.center) to (5.center);
		\draw [style=geometric, in=270, out=-105, looseness=0.75] (6.center) to (7.center);
		\draw [style=geometric, in=90, out=90, looseness=0.75] (8.center) to (9.center);
		\draw [style=geometric, in=270, out=-105, looseness=0.75] (10.center) to (11.center);
		\draw [style=geometric, in=90, out=90, looseness=0.75] (12.center) to (13.center);
		\draw [style=dir] (15.center) to (14.center);
		\draw [style=dir] (17.center) to (16.center);
		\draw [style=geometric,line width=1.05pt] (-4.75, -6.20) to (-4.75, -5.80);
		\draw [style=geometric,line width=1.05pt] (0, -6.20) to (0, -5.80);
		\draw [style=geometric,line width=1.05pt] (4.25, -6.20) to (4.25, -5.80);
		\draw [style=dir,line width=1.05pt] (-4.75, -4.45) to (-4.75, -5.62);
		\draw [style=dir,line width=1.05pt] (0, -4.45) to (0, -5.62);
		\draw [style=dir,line width=1.05pt] (4.25, -4.45) to (4.25, -5.62);
		\draw [in=180, out=-180, looseness=0.75] (49.center) to (48.center);
		\draw [in=180, out=-180, looseness=0.75] (47.center) to (46.center);
		\draw [in=180, out=-180, looseness=0.50] (53.center) to (52.center);
		\draw [in=180, out=-180, looseness=0.50] (51.center) to (50.center);
		\draw [in=180, out=-180, looseness=0.25] (55.center) to (54.center);
	\end{pgfonlayer}
\end{tikzpicture}
}
\end{equation} 
where the arrows identify the levels $t_i$ and the corresponding cuts. Lurie's idea is thus to define the classifying space for composable chains of bordisms of length $k$ by considering the set of all smooth embeddings of $d$-dimensional unoriented manifolds $M \hookra V \times \RR$ together with a choice $\{t_0\leq \cdots \leq t_k\} \subset \RR$ of $k$ points and such that the restricted projection  $\pi_{|M}\maps M \fib \RR$ is proper, a necessary technical condition we will unpack afterwards, and whose set of critical values $\mathsf{Crit}(\pi_{|M})$ is disjoint from $\{t_0\leq \cdots \leq t_k\}$. Importantly, this set can be endowed with a topology which encodes the correct homotopy type --- and although details matter here, we will leave a deeper discussion for the end of this chapter.


Lurie's definition \cite[Def 2.2.9]{Lur09}, despite requiring a lot of legwork to set up properly, offers a perfectly valid construction of a Segal space, which can be further completed to encode the $(\oo,1)$-category of unoriented cobordisms. Despite this, the Segal condition proves hard to check, and it does so for reasons which are unsurprising at this stage: cutting manifolds across a single submanifold, or gluing them if you see it as the reverse operation, is easier to do if you first specify collars. Failure to do so begets issues of coherence and pseudo-isotopy — which, although not intractable, increase the burden of working with such models. This issue can be ameliorated by placing the legwork elsewhere, and this is at the core of the variant introduced in the work of Calaque--Scheimbauer. 

To build collars into the definition, Calaque--Scheimbauer introduce an auxiliary $(\oo,1)$-category of closed intervals $\cI$. Calaque--Scheimbauer denote this category by $\mathrm{Int}$; we use $\cI$ to distinguish it from our interval category $\Int$. As a complete Segal space, its $k$-levels are defined as $(k+1)$-tuples 
\begin{equation}
    \cI_k = \{([a_0,b_0]\leq \cdots \leq [a_k,b_k])\,:\, [a_i,b_i]\subset \RR \}
\end{equation}
topologized via the standard Euclidean topology on $\RR^{2k+2}$. With these intervals, they improve on Lurie's definition by pairing the embedding $M \hookra V\times \RR$ with a choice of a $k$-simplex $(I_0\leq \cdots \leq I_k)\in \cI_k$ such that $\mathsf{Crit}_{|\pi}$ is instead disjoint from $I_0 \cup \cdots \cup I_k$. Visually, we can depict a $k$-simplex as follows 
\begin{equation}
    \scalebox{0.8}{\tikzfig{figures/Ch3/BOR7}{\begin{tikzpicture}
	\begin{pgfonlayer}{nodelayer}
		\node [style=none] (0) at (-9.5, 0) {};
		\node [style=none] (1) at (-4.5, 3) {};
		\node [style=none] (2) at (-4.5, -3) {};
		\node [style=none] (3) at (9.5, 0) {};
		\node [style=none] (4) at (-4.5, 3) {};
		\node [style=none] (5) at (-4.5, -3) {};
		\node [style=none] (6) at (-7, 0.25) {};
		\node [style=none] (7) at (-2, 0.25) {};
		\node [style=none] (8) at (-6.75, -0.35) {};
		\node [style=none] (9) at (-2.5, -0.425) {};
		\node [style=none] (10) at (1.5, 0.25) {};
		\node [style=none] (11) at (6.5, 0.25) {};
		\node [style=none] (12) at (1.75, -0.35) {};
		\node [style=none] (13) at (6, -0.425) {};
		\node [style=none] (14) at (-12, 7) {};
		\node [style=none] (15) at (-12, -6.5) {};
		\node [style=none] (16) at (12, -6) {};
		\node [style=none] (17) at (-12.5, -6) {};
		\node [style=none] (18) at (-13, 7) {$V$};
		\node [style=none] (19) at (12, -6.75) {};
		\node [style=none] (32) at (-4.75, -6) {};
		\node [style=none] (41) at (0, -6) {};
		\node [style=none] (42) at (4.25, -6) {};
		\node [style=none] (46) at (-4.75, -3) {};
		\node [style=none] (47) at (-4.75, -0.85) {};
		\node [style=none] (48) at (-4.75, 0.55) {};
		\node [style=none] (49) at (-4.75, 3) {};
		\node [style=none] (50) at (4.25, -3.65) {};
		\node [style=none] (51) at (4.25, -0.85) {};
		\node [style=none] (52) at (4.25, 0.55) {};
		\node [style=none] (53) at (4.325, 3.625) {};
		\node [style=none] (54) at (0, -3.325) {};
		\node [style=none] (55) at (-0.05, 3.325) {};
		\node [style=none] (56) at (-5.75, -6) {};
		\node [style=none] (57) at (-4, -6) {};
		\node [style=none] (58) at (-1, -6) {};
		\node [style=none] (59) at (1.25, -6) {};
		\node [style=none] (60) at (2.75, -6) {};
		\node [style=none] (61) at (5, -6) {};
		\node [style=none] (62) at (-5.75, 0.475) {};
		\node [style=none] (63) at (-4, 0.525) {};
		\node [style=none] (64) at (-5.875, 2.9) {};
		\node [style=none] (65) at (-3.75, 3.025) {};
		\node [style=none] (66) at (-5.85, -2.925) {};
		\node [style=none] (67) at (-4.1, -3) {};
		\node [style=none] (68) at (-5.925, -0.7) {};
		\node [style=none] (69) at (-3.85, -0.825) {};
		\node [style=none] (70) at (-1.1, -3.225) {};
		\node [style=none] (71) at (1.15, -3.475) {};
		\node [style=none] (72) at (-1, 3.225) {};
		\node [style=none] (73) at (1.075, 3.425) {};
		\node [style=none] (74) at (2.65, -3.6) {};
		\node [style=none] (75) at (5.025, -3.65) {};
		\node [style=none] (76) at (2.75, -0.75) {};
		\node [style=none] (77) at (5.125, -0.75) {};
		\node [style=none] (78) at (2.8, 0.475) {};
		\node [style=none] (79) at (5.15, 0.4) {};
		\node [style=none] (80) at (2.825, 3.6) {};
		\node [style=none] (81) at (5.15, 3.625) {};
		\node [style=none] (82) at (-5.75, -6.75) {$a_0$};
		\node [style=none] (83) at (-4, -6.75) {$b_0$};
		\node [style=none] (84) at (-5.75, -5.75) {};
		\node [style=none] (85) at (-5.75, -6.25) {};
		\node [style=none] (86) at (-4, -5.75) {};
		\node [style=none] (87) at (-4, -6.25) {};
		\node [style=none] (88) at (-1, -5.75) {};
		\node [style=none] (89) at (-1, -6.25) {};
		\node [style=none] (90) at (1.25, -5.75) {};
		\node [style=none] (91) at (1.25, -6.25) {};
		\node [style=none] (92) at (2.75, -5.75) {};
		\node [style=none] (93) at (2.75, -6.25) {};
		\node [style=none] (94) at (5, -5.75) {};
		\node [style=none] (95) at (5, -6.25) {};
		\node [style=none] (96) at (-1, -6.75) {$a_1$};
		\node [style=none] (97) at (1.25, -6.75) {$b_1$};
		\node [style=none] (98) at (2.75, -6.75) {$a_2$};
		\node [style=none] (99) at (5, -6.75) {$b_2$};
	\end{pgfonlayer}
	\begin{pgfonlayer}{edgelayer}
		\draw [style=geometric, in=90, out=-180] (1.center) to (0.center);
		\draw [style=geometric, in=180, out=-90] (0.center) to (2.center);
		\draw [style=geometric, in=90, out=0] (4.center) to (3.center);
		\draw [style=geometric, in=0, out=-90] (3.center) to (5.center);
		\draw [style=geometric, in=270, out=-105, looseness=0.75] (6.center) to (7.center);
		\draw [style=geometric, in=90, out=90, looseness=0.75] (8.center) to (9.center);
		\draw [style=geometric, in=270, out=-105, looseness=0.75] (10.center) to (11.center);
		\draw [style=geometric, in=90, out=90, looseness=0.75] (12.center) to (13.center);
		\draw [style=dir] (15.center) to (14.center);
		\draw [style=dir] (17.center) to (16.center);
		\draw [in=180, out=-180, looseness=0.50] (49.center) to (48.center);
		\draw [in=180, out=-180, looseness=0.75] (47.center) to (46.center);
		\draw [in=180, out=-180, looseness=0.50] (53.center) to (52.center);
		\draw [in=180, out=-180, looseness=0.50] (51.center) to (50.center);
		\draw [in=180, out=-180, looseness=0.25] (55.center) to (54.center);
		\draw [style={open_int}] (56.center) to (57.center);
		\draw [style={open_int}] (59.center) to (58.center);
		\draw [style={open_int}] (60.center) to (61.center);
		\draw [style=red shade] (65.center)
			 to [in=30, out=-180, looseness=0.25] (64.center)
			 to [in=-165, out=-165, looseness=0.50] (62.center)
			 to [in=165, out=15, looseness=0.50] (63.center)
			 to [in=180, out=180, looseness=0.50] cycle;
		\draw [style=red shade] (69.center)
			 to [in=-30, out=-180, looseness=0.50] (68.center)
			 to [in=-165, out=-165, looseness=0.50] (66.center)
			 to [in=180, out=0, looseness=0.50] (67.center)
			 to [in=-180, out=-180, looseness=0.50] cycle;
		\draw [style=red shade] (71.center)
			 to [in=180, out=180, looseness=0.25] (73.center)
			 to [in=0, out=-180, looseness=0.25] (72.center)
			 to [in=-180, out=-180, looseness=0.25] (70.center)
			 to [in=165, out=0, looseness=0.50] cycle;
		\draw [style=red shade] (75.center)
			 to [in=-165, out=180, looseness=0.50] (77.center)
			 to [in=-15, out=-165, looseness=0.50] (76.center)
			 to [in=-165, out=-165, looseness=0.50] (74.center)
			 to [in=180, out=0, looseness=0.50] cycle;
		\draw [style=red shade] (80.center)
			 to [in=-165, out=-165, looseness=0.50] (78.center)
			 to [in=165, out=15, looseness=0.50] (79.center)
			 to [in=180, out=-180, looseness=0.50] (81.center)
			 to [in=0, out=165, looseness=0.25] cycle;
		\draw [style=red shade, in=135, out=45, loop] (56.center) to ();
		\draw [style=red shade, in=135, out=45, loop] (56.center) to ();
		\draw [style=red shade, in=135, out=45, loop] (56.center) to ();
		\draw [style=red shade, in=135, out=45, loop] (85.center) to ();
		\draw [style={open_int}] (84.center) to (85.center);
		\draw [style=red shade, in=135, out=45, loop] (87.center) to ();
		\draw [style={open_int}] (86.center) to (87.center);
		\draw [style=red shade, in=135, out=45, loop] (89.center) to ();
		\draw [style={open_int}] (88.center) to (89.center);
		\draw [style=red shade, in=135, out=45, loop] (91.center) to ();
		\draw [style={open_int}] (90.center) to (91.center);
		\draw [style=red shade, in=135, out=45, loop] (93.center) to ();
		\draw [style={open_int}] (92.center) to (93.center);
		\draw [style=red shade, in=135, out=45, loop] (95.center) to ();
		\draw [style={open_int}] (94.center) to (95.center);
	\end{pgfonlayer}
\end{tikzpicture}
}}
\end{equation}
where the red shaded regions in $M$ are manifestly the pre-images of $\pi_{|M}$ over the intervals $I_i$, and correspond to the collars along which the bordism decomposes. This definition has a slightly larger amount of overhead than Lurie's, but it offers a clearer way of checking the Segal condition and it fixes a failure of essential constancy. 

Our definition tries to reduce the extra machinery in Calaque--Scheimbauer's interval $(\oo,1)$-category $\cI$ by, once again, placing it elsewhere. We use interval Segal spaces, and with this model intervals are part and parcel of the definition of $(\oo,1)$-categories, rather than an auxiliary tool which merely supports the construction of bordism categories. This approach will prove especially useful when we decide to construct the $(\oo,n)$-category of conically smooth stratified cobordisms --- offering almost a machine for churning out higher categories of topological objects embedded in Euclidean space.

In our construction we no longer talk about $k$-simplices as direct encoders of chains of composable morphisms, instead we use a disjoint union of open intervals $I \subset [0,1]$ containing the endpoints. An $I$-level element will be, similarly to the above constructions, an embedding $M \hookra V\times [0,1]$, which we depict as
\begin{equation}
    \scalebox{0.8}{\tikzfig{figures/Ch3/BOR8}{\begin{tikzpicture}
	\begin{pgfonlayer}{nodelayer}
		\node [style=none] (1) at (-4.5, 3) {};
		\node [style=none] (2) at (-4.5, -3) {};
		\node [style=none] (3) at (9.5, 0) {};
		\node [style=none] (4) at (-4.5, 3) {};
		\node [style=none] (5) at (-4.5, -3) {};
		\node [style=none] (7) at (-2, 0.25) {};
		\node [style=none] (10) at (1.5, 0.25) {};
		\node [style=none] (11) at (6.5, 0.25) {};
		\node [style=none] (12) at (1.75, -0.35) {};
		\node [style=none] (13) at (6, -0.425) {};
		\node [style=none] (14) at (-12, 7) {};
		\node [style=none] (15) at (-12, -6.5) {};
		\node [style=none] (16) at (13.75, -6) {};
		\node [style=none] (17) at (-12.5, -6) {};
		\node [style=none] (18) at (-13, 7) {$V$};
		\node [style=none] (41) at (0, -6) {};
		\node [style=none] (42) at (4.25, -6) {};
		\node [style=none] (46) at (-4.75, -3) {};
		\node [style=none] (47) at (-4.75, -0.85) {};
		\node [style=none] (48) at (-4.75, 0.55) {};
		\node [style=none] (49) at (-4.75, 3) {};
		\node [style=none] (50) at (4.25, -3.65) {};
		\node [style=none] (51) at (4.25, -0.85) {};
		\node [style=none] (52) at (4.25, 0.55) {};
		\node [style=none] (53) at (4.325, 3.625) {};
		\node [style=none] (55) at (-0.05, 3.325) {};
		\node [style=none] (57) at (-4, -6) {};
		\node [style=none] (58) at (-1, -6) {};
		\node [style=none] (59) at (1.25, -6) {};
		\node [style=none] (60) at (2.75, -6) {};
		\node [style=none] (61) at (5, -6) {};
		\node [style=none] (81) at (5, 3.625) {};
		\node [style=none] (82) at (-11, -6) {};
		\node [style=none] (83) at (-12, -6) {};
		\node [style=none] (84) at (-10.25, -6) {};
		\node [style=none] (85) at (-6, -6) {};
		\node [style=none] (86) at (10, -6) {};
		\node [style=none] (87) at (11.25, -6) {};
		\node [style=none] (88) at (13, -6) {};
		\node [style=none] (90) at (-5, 0.5) {};
		\node [style=none] (92) at (-7.25, 0.25) {};
		\node [style=none] (93) at (-9.5, -10) {};
		\node [style=none] (100) at (-9.5, 0) {};
		\node [style=none] (101) at (-6.95, -0.225) {};
		\node [style=none] (102) at (-8.175, -2.225) {};
		\node [style=none] (103) at (-7.95, 2.325) {};
		\node [style=none] (104) at (-6, 2.925) {};
		\node [style=none] (105) at (-6, 0.4) {};
		\node [style=none] (106) at (-6, -0.725) {};
		\node [style=none] (107) at (-6, -2.875) {};
		\node [style=none] (108) at (9.5, 0) {};
		\node [style=none] (109) at (6.025, -0.4) {};
		\node [style=none] (110) at (7.975, 2.95) {};
		\node [style=none] (111) at (7.825, -3.075) {};
		\node [style=none] (112) at (5, -3.65) {};
		\node [style=none] (113) at (5, -0.75) {};
		\node [style=none] (114) at (5, 0.425) {};
		\node [style=none] (115) at (5, 3.625) {};
		\node [style=none] (116) at (-2.5, -0.4) {};
		\node [style=none] (118) at (-1, 3.25) {};
		\node [style=none] (119) at (-2.5, 3.1) {};
		\node [style=none] (120) at (-4, 3) {};
		\node [style=none] (121) at (-4, 0.475) {};
		\node [style=none] (122) at (-4, -0.875) {};
		\node [style=none] (123) at (-4, -3) {};
		\node [style=none] (124) at (-2.25, -3.125) {};
		\node [style=none] (125) at (-1, -3.25) {};
		\node [style=none] (126) at (-0.05, -3.3) {};
		\node [style=none] (127) at (1.75, -0.35) {};
		\node [style=none] (128) at (1.25, 3.45) {};
		\node [style=none] (129) at (2, 3.525) {};
		\node [style=none] (130) at (2.825, 3.575) {};
		\node [style=none] (131) at (2.775, 0.475) {};
		\node [style=none] (132) at (2.75, -0.775) {};
		\node [style=none] (133) at (2.8, -3.575) {};
		\node [style=none] (134) at (2.05, -3.525) {};
		\node [style=none] (135) at (1.25, -3.45) {};
	\end{pgfonlayer}
	\begin{pgfonlayer}{edgelayer}
		\draw [style=geometric, in=0, out=-90] (3.center) to (5.center);
		\draw [style=geometric, in=90, out=90, looseness=0.75] (12.center) to (13.center);
		\draw [style=dir] (15.center) to (14.center);
		\draw [style=dir] (17.center) to (16.center);
		\draw [in=180, out=-180, looseness=0.50] (49.center) to (48.center);
		\draw [in=180, out=-180, looseness=0.75] (47.center) to (46.center);
		\draw [in=180, out=-180, looseness=0.50] (53.center) to (52.center);
		\draw [in=180, out=-180, looseness=0.50] (51.center) to (50.center);
		\draw [style={open_int}] (57.center) to (58.center);
		\draw [style={open_int}] (59.center) to (60.center);
		\draw [style={open_int}] (83.center) to (82.center);
		\draw [style={open_int}] (84.center) to (85.center);
		\draw [style={open_int}] (61.center) to (86.center);
		\draw [style={open_int}] (87.center) to (88.center);
		\draw [style=geometric, in=180, out=90] (100.center) to (4.center);
		\draw [style=geometric, in=180, out=-90] (100.center) to (5.center);
		\draw [style=geometric, in=180, out=45] (101.center) to (48.center);
		\draw [style=red shade] (100.center)
			 to [in=-150, out=90] (103.center)
			 to [in=-180, out=30, looseness=0.50] (104.center)
			 to [bend left=270, looseness=0.50] (105.center)
			 to [in=45, out=-165] (101.center)
			 to [in=165, out=-45] (106.center)
			 to [in=-165, out=165, looseness=0.50] (107.center)
			 to [in=-30, out=180, looseness=0.75] (102.center)
			 to [in=-90, out=150] cycle;
		\draw [style=red shade] (108.center)
			 to [in=30, out=-90] (111.center)
			 to [in=0, out=-150, looseness=0.50] (112.center)
			 to [bend left=90, looseness=0.50] (113.center)
			 to [in=-150, out=15] (109.center)
			 to [in=-15, out=105] (114.center)
			 to [bend left=90, looseness=0.50] (115.center)
			 to [in=150, out=0, looseness=0.75] (110.center)
			 to [in=90, out=-30] cycle;
		\draw [style=geometric, in=270, out=-105, looseness=0.75] (10.center) to (11.center);
		\draw [style=geometric, in=255, out=-75, looseness=0.75] (92.center) to (7.center);
		\draw [style=geometric, in=120, out=0, looseness=0.75] (48.center) to (116.center);
		\draw [style=red shade] (125.center)
			 to [in=360, out=165, looseness=0.75] (124.center)
			 to (123.center)
			 to [bend left=90, looseness=0.50] (122.center)
			 to [in=-150, out=0, looseness=0.50] (116.center)
			 to [in=-15, out=135, looseness=1.25] (121.center)
			 to [bend left=90, looseness=0.50] (120.center)
			 to [in=-180, out=0] (119.center)
			 to [in=-165, out=0] (118.center)
			 to [in=180, out=-165, looseness=0.25] cycle;
		\draw [style=geometric, in=180, out=-180, looseness=0.25] (55.center) to (126.center);
		\draw [style=geometric, in=90, out=0] (4.center) to (3.center);
		\draw [style=red shade] (135.center)
			 to [in=-180, out=0, looseness=0.75] (134.center)
			 to (133.center)
			 to [bend left=90, looseness=0.50] (132.center)
			 to [in=-45, out=165, looseness=0.75] (127.center)
			 to [in=-165, out=75] (131.center)
			 to [bend left=90, looseness=0.50] (130.center)
			 to [in=0, out=-180] (129.center)
			 to [in=0, out=180, looseness=0.50] (128.center)
			 to [bend right=90, looseness=0.25] cycle;
	\end{pgfonlayer}
\end{tikzpicture}
}}
\end{equation}
where the restricted projection $\pi_{|M}\maps M \fib [0,1]$ must be proper and such that $\mathsf{Crit}(\pi_{|M})$ is disjoint from $[0,1]\setminus I$. Note that the red shaded regions of $M$ correspond to the pre-image of $I$ under the restricted projection $\pi_{|M}$, and they contain the critical points of the map --- it is over these regions that the cobordism will undergo meaningful changes in topology as it may contain any finite number of critical points with respect to the projection --- on the other hand, in its complementary regions the property outlined above implies that the embedding must be diffeomorphic to a cylinder of a $(d-1)$-dimensional manifold. We will explore this model and its subtle distinctions to the ones discussed above in the next sections.





\subsection{Sheaves on manifolds}

One of the distinctions in our methodology when approaching higher categories of cobordisms, as compared to Lurie \cite{Lur09} and Calaque--Scheimbauer \cite{CS19}, has to do with how we set up the spatial structure. Before any technical detail it is useful to record the aim of this section via the slogan: \emph{offering tools to convert simplicial homotopy theory into smooth homotopy theory}. To properly set up the importance of this issue, we should explain what fails in aforementioned works.

In the work of Calaque--Scheimbauer, the spatial structure of the $n$-fold Segal spaces is constructed with recourse to the smooth singular complex. More concretely, Lemma 5.7 of \cite{CS19} states that the smooth $l$-simplices of a certain topological space can be neatly described as a pair of an embedding $(M \hookra V \times \RR)$ and a sequence of intervals $I_0 \leq \cdots \leq I_k$ satisfying three conditions. From \cite[Corollary 4.47]{CW}, it is understood that a manifold $M$ with boundary or corners is not fibrant as a diffeological space; meaning its associated smooth singular complex $\mathsf{Sing}^{\sf sm}(M)$ is not, in general, a Kan complex. Tracking the reasoning in \cite{CS19}, it is unclear how to ensure that the topological space in question has the fibrancy properties required for the Segal-space argument.

To avoid dealing with the cumbersome point-set topologies on spaces of submanifolds, one can instead work with simplicial sets. These, however, are rarely Kan complexes. This failure can be detected by testing whether it is possible to glue manifolds along closed subsets. Such gluing procedures, much like gluing smooth functions, are generally impossible along closed sets. This line of thinking is what motivates the appearance of sheaves on the site of smooth manifolds, and in particular the framework developed by Berwick-Evans, Boavida de Brito, and Pavlov. Their theory provides precisely the homotopical tools needed to perform such “gluing” operations along closed subsets, up to smooth homotopy.



\subsection{Sites and sheaves}

As noted above, in contrast to Calaque--Scheimbauer we do not model the spatial
structure of our $(\oo,1)$-categories via smooth singular simplices. Instead,
we work with a homotopical theory of ``spaces'' that has superior gluing
properties — namely, ($\oo$-)sheaves on the site of smooth manifolds.  This sheaf-theoretic viewpoint clarifies the necessary gluing arguments. Our main references for these methods are
\cite{BBEP24} and \cite{MW07}. 

\medskip

We begin by recalling the basic sheaf-theoretic formalism. Let $\cC$ be a small category. The presheaf category $\psh(\cC)$ is the free cocompletion of $\cC$: any diagram in $\cC$ that may fail to admit a colimit does admit one in $\psh(\cC)$ after applying the Yoneda embedding
\begin{equation}
  \cC \hookra \psh(\cC).
\end{equation}
In particular, if $\man$ denotes the category whose objects are
finite-dimensional smooth manifolds without boundary and whose morphisms are
smooth maps, then $\man$ does not admit all colimits (e.g. the span
\[
* \longleftarrow S^0 \xra{\nssphere} S^1
\]
has no colimit in $\man$), whereas $\psh(\man)$ does. However, the colimits in
$\psh(\man)$ are generally \emph{not} the ``correct'' geometric ones.

A simple indication of this failure is disjoint union. Given smooth manifolds
$M_1,M_2$, the coproduct $M_1\sqcup M_2$ exists in $\man$ and is given by
disjoint union. A map
\[
S^0 \longra M_1 \sqcup M_2
\]
amounts to choosing either two points in a single component or one point in
each component. On the other hand, in presheaves we have
\[
\Hom_{\psh(\man)}(rS^0,rM_1 \sqcup rM_2)
\simeq \man(S^0,M_1)\,\sqcup\,\man(S^0,M_2),
\]
so a map $rS^0 \to rM_1 \sqcup rM_2$ is equivalent to choosing two points in
$M_1$ or two points in $M_2$. In other words, we are gluing incorrectly: the
geometric context encoded in $\man$ has been lost. Sheaves repair exactly this
loss by forcing certain ``geometric'' cones to become colimits.

\medskip

Abstractly, the geometric content is encoded by specifying which families of
maps $\{c_i\to c\}$ should behave like covers, and then requiring the
presheaf to reconstruct its value on $c$ as the limit of its values on the
$c_i$ and their intersections. This is the notion of a Grothendieck topology.

\begin{definition}\label{ch3:: def: site}
   Let $\cC$ be a category. A \textbf{Grothendieck topology} on $\cC$ assigns
   to every object $c\in \cC$ a collection $\tau(c)$ of families
   $\{f_i\maps c_i \longra c\}$, called \textbf{covering families}, such that:
   \begin{enumerate}
     \item If $f\maps d \longra c$ is an isomorphism, then
     $\{f\maps d \longra c\}$ is a covering family.
     \item For a covering family $\{f_i \maps c_i \longra c\}$ and any map
     $g\maps d \longra c$, the fiber product
     \[
        \begin{tikzcd}[ampersand replacement=\&]
            {d \times_{c}c_i} \&\& d \\
            \\
            {c_i} \&\& c
            \arrow[from=1-1, to=1-3]
            \arrow[from=1-1, to=3-1]
            \arrow["\lrcorner"{anchor=center, pos=0.125}, draw=none, from=1-1, to=3-3]
            \arrow[from=1-3, to=3-3]
            \arrow[from=3-1, to=3-3]
        \end{tikzcd}
     \]
     exists in $\cC$ and the family of projections
     $\{\pi_i \maps d\times_c c_{i} \longra d\}$ is a covering family.
     \item If $\{f_i \maps c_i \longra c\}$ is a covering family and, for each
     $i$, $\{g_{ji}\maps c_{ji} \longra c_{i}\}$ is a covering family, then
     the composites
     \[
        \{f_i \circ g_{ji}\maps c_{ji} \longra c\}
     \]
     form a covering family of $c$.
   \end{enumerate}
   A covering family with a single map $f\maps d \longra c$ is called a
   \textbf{covering map}. A category $\cC$ endowed with a Grothendieck
   topology $\tau$ is called a \textbf{site}.
\end{definition}

We now introduce the only example of a Grothendieck topology we will need.

\begin{definition}\label{ch3:: def: smooth site}
    Let $\man$ denote the category whose objects are all finite-dimensional
    smooth manifolds without boundary and smooth maps thereamong. The site
    structure is given by \emph{open covers}. For $M\in\man$, a
    \textbf{covering family} is a collection of open embeddings
    \[
      \{U_i \hookra M\}_{i\in I}
    \]
    such that
    \[
      M \,=\, \bigcup_{i \in I} U_i,
    \]
    where we conflate each open embedding with its image.
\end{definition}

\begin{definition}\label{ch3:: def: sheaf}
    Let $\cC$ be a site with Grothendieck topology $\tau$. A functor
    $\cF\maps \cC^{\op} \longra \set$ is called a \textbf{presheaf}. A
    presheaf $\cF$ is a \textbf{sheaf} (with respect to $\tau$) if for each
    covering family $\{f_i\maps c_i \longra c\}_{i\in I}$ the dashed arrow in
    \[
      \cF(c) \dashrightarrow \prod_{i\in I}\cF(c_i)
      \rightrightarrows \prod_{i,j\in I}\cF(c_i\times_c c_j)
    \]
    exhibits $\cF(c)$ as the equalizer of the two parallel arrows.
\end{definition}

In what follows we are primarily interested in sheaves on the smooth site $\man$ of Definition~\ref{ch3:: def: smooth site}.
\bigskip

\subsubsection{Concordance after Madsen--Weiss}

We now recall the notion of concordance from Madsen--Weiss as it is presented in \cite{MW07} --- these ideas can be found in earlier papers of Waldhausen \cite{Wa} and Weiss--Williams \cite{MW95}. This is the basic equivalence relation that will later be promoted to a homotopical condition for $\oo$-sheaves.

\begin{definition}[Concordance]\label{ch3:: def: concordance}
  {\cite[Definition~2.1.3]{MW07}}
  Let $\cF\maps \man^{\op} \longra \set$ be a presheaf and let
  $\pi \maps X \times \RR \longra \RR$ be the projection. Two sections
  $\sigma_0,\sigma_1 \in \cF(X)$ are \textbf{concordant} if there exists a
  section $\Sigma \in \cF(X\times\RR)$ such that:
  \begin{enumerate}
    \item There is an open neighborhood
    $X \times (-\oo,0] \subset U_{-}$ with
    $\Sigma_{|U_{-}} = (\pi^* \sigma_{0})_{|U_{-}}$.
    \item There is an open neighborhood
    $X \times [0,+\oo) \subset U_{+}$ with
    $\Sigma_{|U_{+}} = (\pi^* \sigma_{1})_{|U_{+}}$.
  \end{enumerate}
  The section $\Sigma$ is called a \textbf{concordance} from $\sigma_0$ to
  $\sigma_1$, and we write $\sigma_0 \sim_c \sigma_1$.
\end{definition}

\begin{example}
  Let $N$ be a smooth manifold and let
  \[
    \cF = C^{\oo}(-,N)\maps \man^{\op} \longra \set
  \]
  be the presheaf $M \mapsto C^{\oo}(M,N)$. Given $f,g\in C^{\oo}(M,N)$, a
  concordance $H\in C^{\oo}(M\times\RR,N)$ is precisely a smooth homotopy
  $H\maps M\times\RR\to N$ from $f$ to $g$ that is eventually constant in
  each end, i.e.\ $H(x,t)=f(x)$ for $t<0$ and $H(x,t)=g(x)$ for $t>0$.
\end{example}

\begin{lemma}\label{ch3:: lem: concordance is equiv}
    Let $\cF\maps \man^\op \longra \set$ be a sheaf. Then the relation
    $\sim_c$ on $\cF(X)$ is an equivalence relation.
\end{lemma}

\begin{proof}
  Reflexivity is witnessed by $\pi^*\sigma\in\cF(X\times\RR)$ for any
  $\sigma\in\cF(X)$. Symmetry is obtained by precomposing a concordance
  $\Sigma$ with $-\id_{\RR}$. For transitivity, suppose
  $\sigma_0\sim_c\sigma_1$ and $\sigma_1\sim_c\sigma_2$ with concordances
  $\Sigma_{01},\Sigma_{12}\in\cF(X\times\RR)$. Choosing a diffeomorphism
  $\phi\maps\RR\to\RR$ that sends an open neighborhood of $+\infty$ to all of
  $\RR$ and reparameterizing $\Sigma_{12}$ along $\phi$, one can glue the two
  concordances along an open neighborhood of $X\times\{0\}$ using the sheaf
  property, and then reparameterize again to obtain a concordance from
  $\sigma_0$ to $\sigma_2$. Details follow the argument in \cite[Lemma~3.1]{BBEP24}.
\end{proof}

\begin{remark}\label{ch3:: rmk: MW-representability}
  Given a sheaf $\cF\maps\man^\op\to\set$ and $X\in\man$, consider the set
  \[
    \cF[X] \;\overset{\sf def}{=}\; \cF(X)/\sim_c.
  \]
  This assignment is functorial in $X$ and defines a presheaf
  $\cF[\bullet]\maps\man^\op\to\set$. It is rarely a sheaf again. Madsen--Weiss
  show that this presheaf is nevertheless representable in a weak sense:
  there exists a space $|\cF|$ such that for every manifold $X$ there is a
  natural bijection
  \[
    \pi_0\Map(X,|\cF|)\;\simeq\;\cF[X].
  \]
  Moreover, $|\cF|$ can be constructed as the geometric realization of the
  simplicial set $[n]\mapsto \cF(\Delta^n_e)$, where $\Delta^n_e$ is the
  extended simplex in the sense of \cite[Definition~2.4.2]{MW07}, and this
  representability extends to relative concordance classes and higher
  homotopy groups.
\end{remark}

Madsen--Weiss work with set-valued sheaves. In our application we need the
fully homotopical version developed by Berwick-Evans--Boavida de Brito--Pavlov
\cite{BBEP24}, where one replaces sheaves of sets by $\oo$-sheaves of spaces.
\bigskip

\subsubsection{\texorpdfstring{$\oo$}{infty}-sheaves and the concordance resolution}

We now upgrade the previous discussion to $\oo$-sheaves. Following
\cite{BBEP24}, by a \emph{space} we mean a simplicial set, and by a
\emph{simplicial space} a simplicial object in spaces.

\begin{definition}\label{ch3:: def: infinity-sheaf}
  {\cite[Definition~2.15]{BBEP24}}
  A presheaf $\cF\maps \man^{\op}\to\sset$ is an \textbf{$\oo$-sheaf} if for every manifold $M$ and every open cover \[\{U_i\hookra M\}_{i\in I},\] the  canonical map
  \begin{equation}
    \cF(M)\;\longra\;
    \hoLimit[\bDelta]\Big(
      \prod_{i_0}\cF(U_{i_0})
      \doubleleftarrow
      \prod_{i_0,i_1}\cF(U_{i_0}\cap U_{i_1})
      \tripleleftarrow\cdots
    \Big)
   \end{equation}
  is a weak equivalence of spaces.
\end{definition}


\begin{remark}
    The homotopy limit in the definition of an $\oo$-sheaf for a given open cover $\cU = \{U_\alpha \hookra M\}$ of $M$, is taken over the cosimplicial space 
    \begin{equation}
        \bDelta \xlonghookra{\,\,\,\,\, (\check{\cU}_{\bullet})^{\op}\,\,\,\,\,\,} \man^\op \xlonghookra{\,\, \cF \,\,} \sset,
    \end{equation} 
    where $\check{\cU}_{\bullet}$ is the Cech complex
    \begin{align}
        \check{\cU}_{\bullet} \maps \bDelta^{\op} &\longra \man \\
        [n] &\longmapsto \coprod_{i_{0},\cdots,i_{n}} U_{i_{0}}\cap U_{i_{1}} \cap \cdots \cap U_{i_{n}}.
    \end{align}
\end{remark}

\begin{theorem}[Binary descent criterion for \(\oo\)-sheaves]\label{ch3:: thm: BBEP-sheaf-criterion}
{\cite[Theorem~5.1]{BBEP24}}
A presheaf \(\cF\colon \man^{\op}\to\sset\) is an \(\oo\)-sheaf (i.e.\ satisfies descent for open covers)
if and only if \(\cF(\emptyset)\simeq *\) and the following two conditions hold:
\begin{enumerate}
  \item \label{binary cover} \textbf{Binary open covers}: For every manifold \(M\) and every open cover \(M=U\cup V\),
  the commutative square
  \begin{equation}
    \begin{tikzcd}[ampersand replacement=\&]
        {\cF(M)} \&\& {\cF(V)} \\
        \\
        {\cF(U)} \&\& {\cF(U\cap V)}
        \arrow[from=1-1, to=1-3]
        \arrow[from=1-1, to=3-1]
        \arrow[from=1-3, to=3-3]
        \arrow[from=3-1, to=3-3]
    \end{tikzcd}
  \end{equation}
  is a homotopy pullback square in \(\sset\).
  \item \label{disj union} \textbf{Disjoint unions}: If 
  \begin{equation}
    M \cong \bigsqcup_{i\in I} U_i
  \end{equation} is a (possibly infinite) disjoint union of
  submanifolds \(U_i\), then the canonical map
  \begin{equation}
    \cF(M)\;\longra\;\prod\nolimits^{h}_{i\in I} \cF(U_i)
  \end{equation}
  is a weak equivalence.
\end{enumerate}
\end{theorem}


Set-valued sheaves and groupoid-valued stacks are precisely those presheaves that are $\oo$-sheaves with values in discrete spaces and groupoids, respectively.

\begin{definition}\label{ch3:: def: concordance-invariant}
  {\cite[Definition~2.4]{BBEP24}}
  A presheaf $\cF\maps\man^{\op}\to\sset$ is \textbf{concordance-invariant} if
  for every manifold $M$ the map
  \begin{equation}
    \cF(M)\longra \cF(M\times\RR)
  \end{equation}
  induced by the projection $M\times\RR\to M$ is a weak equivalence.
\end{definition}

Let $\AA^n\subset\RR^{n+1}$ denote the smooth extended $n$-simplex
\begin{equation}
  \AA^n
  \;=\;
  \{(x_0,\dots,x_n)\in\RR^{n+1}\mid x_0+\cdots+x_n=1\},
\end{equation}
which defines a cosimplicial object $\AA^\bullet\maps\bDelta\to\man$. This cosimplicial object endows $\man$ with an enrichment in spaces by declaring 
\begin{equation}
    \Hom_{\man}(M,N)_{\bullet} \; \overset{\sf def}{=} \; \Hom_{\man}(M \times \AA^{\bullet}, N).
\end{equation}
A presheaf $\cF\maps \man^{\op} \longra \sset$ that can be enriched is therefore one which sends smooth homotopies to simplicial homotopies, and smooth homotopy equivalences to simplicial homotopy equivalences. In particular, an enriched presheaf is always concordance-invariant.

Informally, this construction turns a moduli sheaf into a space by retaining
its smooth families: a family over $\AA^1$ gives a concordance, while families
over higher affine simplices encode the corresponding higher homotopies.
For the bordism sheaves below, these are families of embedded bordisms, with
their cutting regions retained throughout the family.

\begin{definition}\label{ch3:: def: concordance-resolution}
  {\cite[Definition~2.5]{BBEP24}}
  Let $\cF\maps\man^{\op}\to\sset$ be a presheaf. Its \textbf{concordance
  resolution} is the presheaf
  \begin{equation}
    \cF(-\times \AA^{\bullet})\maps \man^{\op} \longra \fun(\bDelta^{\op},\sset).
  \end{equation}
  We denote the post-composition of $\cF(-\times \AA^\bullet)$ with the homotopy colimit (of simplicial spaces) as
  \begin{equation}
    \mathbf{B}\cF\maps\man^{\op}\longra\sset,\qquad
    \mathbf{B}\cF(M)\;=\;
    \left\vert\,[k]\longmapsto\cF(M\times \AA^k)\,\right\vert,
  \end{equation}
   We use the diagonal of the associated bisimplicial set as a model here.
\end{definition}


\begin{remark}
    The concordance resolution $\cF(-\times \AA^\bullet)$ lifts to an enriched functor defined on the hom-sets by 
    \begin{equation}
        \Hom_{\man}(M \times \AA^{n}, N) \longra \Hom_{\fun(\bDelta^\op,\sset)}(\cF(N \times \AA^{\bullet})\times \bDelta[n], \cF(M \times \AA^{\bullet})),
    \end{equation}
    for all $n\geq 0$. Post-composing this with the map to its homotopy colimit functor translates to a simplicial map
    \begin{equation}
        \Hom_{\man}(M \times \AA^{n}, N) \longra \Hom_{\sset}\left(\bB\cF(N)\times \Delta^n, \bB\cF(M )\right).
    \end{equation}
    Concluding that $\bB\cF$ is both enriched in spaces, and therefore concordance-invariant — hence the naming choice.
\end{remark}

Berwick-Evans, Boavida de Brito, and Pavlov prove that if $\cF$ is already an $\oo$-sheaf, then $\mathbf{B}\cF$ remains an $\oo$-sheaf and in fact represents concordance classes of sections \cite{BBEP24}.

\begin{theorem}\label{ch3:: thm: BBEP-classifying}
  {\cite[Theorems~1.1 and~1.2]{BBEP24}}
  Let $\cF$ be an $\oo$-sheaf on $\man$. Define the \emph{classifying space} of $\cF$ to be a Kan fibrant replacement $B\cF$ of the simplicial set $\mathbf{B}\cF(*)$. Then:
  \begin{enumerate}
    \item For every manifold $M$ there is a natural weak equivalence
      \begin{equation}
        (\mathbf{B}\cF)(M)\;\simeq\;
        \Map^{h}(\Sing M, B\cF),
      \end{equation}
      where $\Sing M$ is the singular complex and $\Map^{h}$ denotes the
      derived mapping space.
    \item In particular, $\mathbf{B}\cF$ is an $\oo$-sheaf.
    \item The canonical map
      \[
        \pi_0(\mathbf{B}\cF(M))\;\longrightarrow\;
        \pi_0\big(\cF(M)\big)/\sim_c
      \]
      is a bijection, so $B\cF$ represents concordance classes of sections of
      $\cF$ in the same way that $|\cF|$ represents concordance classes for
      set-valued sheaves in \cite{MW07}. Moreover, for sheaves of sets one
      obtains $\pi_n(B\cF,x)\cong \pi_n$ of the Madsen--Weiss representing
      space.
  \end{enumerate}
\end{theorem}



\subsubsection{Weak Kan fibrations and realization fibrations}

The proof of Theorem~\ref{ch3:: thm: BBEP-classifying} relies on a careful
comparison between the homotopy limit condition defining $\oo$-sheaves and
the homotopy colimits appearing in the concordance resolution. The technical
bridge is the notion of a \emph{weak Kan fibration} of simplicial spaces.

We briefly recall the key definitions from \cite[Section~3]{BBEP24}.

\begin{definition}\label{ch3:: def: weak-RLP}\cite[Definition~3.3]{BBEP24}
  Let $f\maps X\to Y$ be a Reedy fibration between Reedy fibrant simplicial
  spaces, and let $i\maps A\to B$ be a map of simplicial spaces. We say that
  $f$ has the \textbf{weak right lifting property} with respect to $i$ if for
  every commutative square
  \begin{equation}
       \begin{tikzcd}[ampersand replacement=\&]
            A \&\& X \\
            \\
            B \&\& Y
            \arrow[""{name=0, anchor=center, inner sep=0}, "\beta", from=1-1, to=1-3]
            \arrow["i"', from=1-1, to=3-1]
            \arrow["f", from=1-3, to=3-3]
            \arrow["{\tilde{\alpha}}"{description}, dashed, from=3-1, to=1-3]
            \arrow["\alpha"', from=3-1, to=3-3]
            \arrow["H"{description}, color={rgb,255:red,214;green,92;blue,92}, between={0.3}{0.7}, Rightarrow, from=0, to=3-1]
        \end{tikzcd}
    \end{equation}
  there exists a map $\tilde\alpha\colon B\to X$ and a \textbf{vertical
  homotopy} \[H\colon A\times\Delta[1]\to X\] satisfying the following conditions:
      \begin{enumerate}
        \item $A\times \Delta[\{0\}] \hookra A \times \Delta[1] \xra{H} X= \beta$.
        \item $A\times \Delta[\{1\}] \hookra A \times \Delta[1] \xra{H} X = \widetilde{\alpha}\circ i$.
        \item The following diagram commutes:
        \begin{equation}
        \begin{tikzcd}[ampersand replacement=\&]
            {A\times \Delta[1]} \&\& X \&\& Y \\
            \\
            A \&\& B
            \arrow["H", from=1-1, to=1-3]
            \arrow["\pi"', from=1-1, to=3-1]
            \arrow["f", from=1-3, to=1-5]
            \arrow["i"', hook, from=3-1, to=3-3]
            \arrow["\alpha"', from=3-3, to=1-5]
            \arrow["\tilde{\alpha}", dashed, from=3-3, to=1-3]
        \end{tikzcd}
        \end{equation}
    \end{enumerate}
\end{definition}

Informally, weak lifting translates to the ability of solving lifting problems for $i$ up to a specified homotopy that is \emph{vertical} over $Y$. For a more illustrative explication of the motivation behind this definition see \cite[Remark~3.4]{BBEP24}

\begin{definition}\label{ch3:: def: weak-kan-fibration}
  {\cite[Definition~3.8]{BBEP24}}
  A Reedy fibration $f\maps X\to Y$ of simplicial spaces is a
  \textbf{weak Kan fibration} if it has the weak right lifting property with
  respect to the maps
  \[
    h_{i,n,k}\colon
    \mathsf{Sd}^{i}\big(\Lambda_k[n]\big)\longrightarrow
    \mathsf{Sd}^{i}\big(\Delta[n]\big)
  \]
  for all $i\geq 0$, $n\geq1$ and $0\leq k\leq n$, where
  $\mathsf{Sd}^{i}$ denotes the $i$-fold subdivision functor on simplicial
  spaces.
\end{definition}

This notion interpolates between genuine Kan fibrations and the geometric fibrations needed for concordance resolutions. The key property is that weak Kan fibrations control the interaction between homotopy pullbacks and geometric realization.

\begin{definition}\label{ch3:: def: realization-fibration}
  A map $f\maps X\to Y$ of simplicial spaces is a
  \textbf{realization fibration} if for every $Z\to Y$ the induced map
  \begin{equation}
    \big|X\times_Y^{h} Z\big|
      \longrightarrow
    |X|\times_{|Y|}^{h}|Z|
  \end{equation}
  is a weak equivalence of spaces.
\end{definition}

\begin{theorem}\label{ch3:: thm: weak-kan-realization}
  {\cite[Theorem~3.17]{BBEP24}}
  A weak Kan fibration is a realization fibration.
\end{theorem}

In particular, geometric realization commutes with homotopy pullbacks along
weak Kan fibrations. This is the mechanism allowing BBEP to commute the
homotopy limits in the $\oo$-sheaf condition with the homotopy colimits
appearing in the concordance resolution $M\longmapsto\big|[k]\mapsto\cF(M\times \AA^k)\big|$, and thus to prove Theorem~\ref{ch3:: thm: BBEP-classifying}.

\begin{example}
    For exposition purposes, let us unpack the example given in \cite[Example~3.18]{BBEP24}, since it highlights the need for the lifts against all simplicial subdivisions when defining weak Kan fibrations. Let $X$ be the union of three non-degenerate 1-simplices and the projection $f\maps X \ra \Delta[1]$ as in the diagram below:
    \begin{equation}
        \begin{tikzcd}[ampersand replacement=\&]
            \bullet \& \bullet \\
            \bullet \& \bullet \\
            \\
            \bullet \& \bullet
            \arrow[from=1-1, to=1-2]
            \arrow[from=1-1, to=2-1]
            \arrow[""{name=0, anchor=center, inner sep=0}, from=2-1, to=2-2]
            \arrow[""{name=1, anchor=center, inner sep=0}, from=4-1, to=4-2]
            \arrow["f", between={0.2}{0.8}, from=0, to=1]
        \end{tikzcd}
    \end{equation} 
    A key couple of observations follow:
    \begin{enumerate}
        \item The map $f$ is not a Kan fibration. The following outer 2-horn (indicated in a separate color)
        \begin{equation}
            \begin{tikzcd}[ampersand replacement=\&]
                \bullet \& \bullet \\
                \bullet \& \bullet
                \arrow[color={rgb,255:red,214;green,92;blue,92}, from=1-1, to=1-2]
                \arrow[color={rgb,255:red,214;green,92;blue,92}, from=1-1, to=2-1]
                \arrow[from=2-1, to=2-2]
            \end{tikzcd}
        \end{equation}  
        over the degenerate 2-simplex, by inspection, does not admit a filler.
        \item Despite the point above, the map $f$ admits weak right lift against $h_0$; on the other hand, it does not admit a weak lift against $h_2$. To understand this example, consider the following map $\mathsf{Sd}^{2}(\Lambda^{2}[0]) \longra X$:
        \begin{equation} 
        \begin{tikzcd}[ampersand replacement=\&]
            \& \bullet \& \bullet \&\&\& \bullet \\
            {} \& \bullet \\
            \& \bullet \& \bullet \& \bullet \& \bullet \& \bullet
            \arrow[from=1-2, to=1-3]
            \arrow[from=1-2, to=2-2]
            \arrow[from=1-6, to=1-3]
            \arrow[from=3-2, to=2-2]
            \arrow[from=3-2, to=3-3]
            \arrow[from=3-4, to=3-3]
            \arrow[from=3-4, to=3-5]
            \arrow[from=3-6, to=3-5]
        \end{tikzcd}
        \end{equation} 
        then any homotopy lifting $\Delta[2] \longra \Delta[1]$ would fail to be vertical.
        \item This example indicates the importance of the weak lifts against repeated subdivisions. Indeed, this map is not a realization fibration: homotopy fiber over $0$ is contractible whilst the homotopy fiber over $1$ is weak homotopy equivalent to a pair of points — therefore $f$ is not a quasi-fibration, thereby not a realization fibration.
    \end{enumerate}
\end{example}


The specific weak Kan fibrations we will use arise from \emph{endpoint
evaluation} for $\oo$-sheaves. Consider the inclusion of the boundary
$\partial \AA^1 = \AA^0\sqcup \AA^0 \hookra \AA^1$, and form the induced map of
simplicial spaces
\begin{equation}
  \cF(\AA^{\bullet}\times \AA^{1}) \longra
  \cF(\AA^{\bullet}\times \partial \AA^{1})
\end{equation}
for an $\oo$-sheaf $\cF\maps\man^{\op}\to\sset$.

\begin{proposition}\label{ch3:: prop: weak-kan-endpoints}\cite[Proposition~4.2]{BBEP24}
  Let $\cF\maps\man^{\op}\to\sset$ be an $\oo$-sheaf. Then the map of
  simplicial spaces
  \begin{equation}
    \cF(\AA^{\bullet}\times \AA^{1}) \longra
    \cF(\AA^{\bullet}\times \partial \AA^{1})
  \end{equation}
  is a weak Kan fibration.
\end{proposition}

This proposition is the crucial glue between the concordance-theoretic framework of Berwick-Evans--Boavida de Brito--Pavlov and the Segal-type constructions which we will entertain in the last section: it allows us to treat endpoint evaluation for concordance data as a type of quasifibration (see \cite{May90}) whose behaviour is well-controlled under realization.

The strategy of using this endpoint weak Kan fibration to transport the concordance-resolution argument of \cite[Theorem~5.1]{BBEP24} into the proof of the interval Segal condition was suggested to the author by Pedro Boavida de Brito.

\medskip

We will apply these general results in the concrete situation where $\cF$ is an $\oo$-sheaf of $d$-dimensional manifolds (and later of bordisms and stratified bordisms). The upshot is that the concordance classes of such geometric data are represented by explicitly described classifying spaces, constructed functorially from $\cF$ via the concordance resolution.

\section{An interval Segal space of \texorpdfstring{$d$}{d}-dimensional bordisms}

We will now construct an interval Segal space whose associated $(\oo,1)$-category is that of $d$-dimensional unoriented bordisms. To construct this, we will make use of the theory of smooth $\oo$-sheaves, i.e.\ $\oo$-sheaves on the site of smooth manifolds, as a conduit for constructing appropriate spaces of smooth manifolds. This adds a distinct flavor to our approach, and clarifies the minor pathologies identified in previous constructions --- as explained in the introduction to this chapter. 

Before proceeding with definitions, let us first sketch the path we will take in this section. We start by defining a general construction
\[
\widetilde{\bord}_{d,V}(-,-)
\]
which is a bifunctor on $\man^{\op}\times \Int$ that, when fixing $T\in \Int$, produces a sheaf of sets on $\man^{\op}$ and captures a type of moduli-space construction for $d$-dimensional smooth manifolds embedded in a fixed finite-dimensional real vector space $V$, fibered over smooth manifolds of varying dimension. It is designed so that it also encompasses our desired classifying spaces of bordisms. Upon taking its concordance resolution and further restricting to $\bDelta^{\op}$, we get a simplicial set $\bord_{d,V}$, or a space, that we will show is the classifying space of $\gamma(T)$-composites of $d$-dimensional bordisms.

Lastly, the most technical part of this section is showing that this composite
construction, initially only a pre-interval space, is in fact an interval Segal
space whose homotopy category is equivalent to the $d$-dimensional bordism
category $\mathsf{Cob}_d$.

\begin{notation}
Given an object $T\in \Int$, we write $T^{\mathsf{inn}}$ for the union of connected open intervals which do not contain the endpoints. For objects $T$ with $0$ connected inner components we have $T^{\mathsf{inn}} = \emptyset$. 
\end{notation}

\begin{notation}
Given an object $I \in \Int$, we call the union
\[
(I^{\mathsf{c}}\cup I^{\mathsf{inn}}) \subset [0,1]
\]
the \textbf{core} of $I$ and denote it by $\mathsf{cor}(I)$. For every $I\in \Int$ the core $\core(I)$ is a connected closed interval of $[0,1]$.
\end{notation}

\begin{definition}\label{ch3:: def: bord_d}
    Let $V$ be a finite-dimensional real vector space. We define a functor 
    \begin{equation}
    \widetilde{\bord}^{B}_{d,V} \maps \mathsf{Man}^{\op}\times \Int \longrightarrow \set    
    \end{equation}
    sending a pair $(M,T)$, where $M\in \man$ and $T\in \Int$, to the set of bounded closed subsets 
    \begin{equation}
    W\subseteq M\times \mathsf{cor}(T) \times V
    \end{equation} 
    satisfying the following three conditions:
    \begin{enumerate}
    \item \label{id: dim} $W$ is a smooth $(d+\mathsf{dim}(M))$-dimensional submanifold with boundary $\partial W$ such that 
    \begin{equation}
        \partial W \subset \partial \big(M \times \mathsf{cor}(T) \times V\big) = M \times \partial (\mathsf{cor}(T)) \times V.   
    \end{equation}
    In other words, the boundary of $W$ is contained in the boundary of the ambient space.
    \item \label{id: categorical} The restricted projection 
    \begin{equation}
        \pi\maps \big(W \cap (M\times T^{\mathsf c} \times V) \big) \hookra  M \times T^{\mathsf c} \times V \fib  M \times T^{\mathsf c}
    \end{equation} 
    is a submersion. 
    \item \label{id: spatial} The projection 
    \begin{equation}
    \pi_W\maps W \hookra M\times \mathsf{cor}(T) \times V \fib M 
    \end{equation}
    is a submersion.
    \end{enumerate} 

    \noindent The functor $\widetilde{\bord}^{B}_{d,V}$ sends:
    \begin{itemize}
    \item[($\man$):] A morphism 
    \[
    (f,\mathsf{id}_{T})\maps (N,T) \longra (M,T)
    \]
    induced by a smooth map $f\maps N \longra M$ to the pullback submanifold 
    \[
    (f \times \id)^{*} W \subseteq N \times \mathsf{cor}(T) \times V,
    \]
    where $W\subseteq M\times \mathsf{cor}(T) \times V$ is a submanifold with boundary as above, and $(f\times \id)^\ast W$ is defined by the pullback diagram
    \begin{equation}
        \begin{tikzcd}[ampersand replacement=\&]
            {(f\times \id)^\ast W} \& {W} \\
            \\
            N\times \mathsf{cor}(T) \times V \& M\times \mathsf{cor}(T) \times V,
            \arrow[from=1-1, to=1-2]
            \arrow[hook, "(f\times \id)^\ast \pi_{|W}"',from=1-1, to=3-1]
            \arrow["\lrcorner"{anchor=center, pos=0.125}, draw=none, from=1-1, to=3-2]
            \arrow[hook, from=1-2, to=3-2]
            \arrow["f\times \id"', from=3-1, to=3-2]
        \end{tikzcd}
    \end{equation}
    which is a $\big(d + \mathsf{dim}(N)\big)$-dimensional submanifold with boundary
    \begin{equation}
        (f\times \id)^* W \subseteq N \times \mathsf{cor}(T) \times V,
    \end{equation}
    since the pullback preserves the fibers over the boundary $\partial (N \times \mathsf{cor}(T) \times V)$.

    \item[($\Int$):] A morphism $(\id_M, T\subset S)$ is defined by \emph{forgetting} connected components and restricting the submanifold with boundary 
    \begin{equation}
        W \subseteq M \times \mathsf{cor}(S) \times V
    \end{equation}
    along the inclusion $M \times \mathsf{cor}(T)\times V \hookra M \times \mathsf{cor}(S)\times V$. We denote the restriction by $W_{|T}$, which comes equipped with a smooth embedding 
    \begin{equation}
        W_{|T} \hookra W.
    \end{equation}
    \end{itemize}  
\end{definition}

\begin{remark}
    Two points are worth recording about this definition. First, regarding well-definedness of the functor, the pullback $(f\times \id)^*W$ admits a smooth $(d+\mathsf{dim}(N))$-dimensional embedded submanifold structure (with boundary) in $N\times \mathsf{cor}(T) \times V$ because, in the following diagram
    \begin{equation}
        \begin{tikzcd}[ampersand replacement=\&]
            {(f\times \id)^\ast W} \&\& W \\
            \\
            {N\times \mathsf{cor}(T) \times \RR^{\oo}} \&\& {M\times \mathsf{cor}(T) \times \RR^{\oo}} \\
            \\
            N \&\& M
            \arrow[from=1-1, to=1-3]
            \arrow[hook, from=1-1, to=3-1]
            \arrow["\lrcorner"{anchor=center, pos=0.125}, draw=none, from=1-1, to=3-3]
            \arrow[hook, from=1-3, to=3-3]
            \arrow["{f\times \id}"', from=3-1, to=3-3]
            \arrow[two heads, from=3-1, to=5-1]
            \arrow["\lrcorner"{anchor=center, pos=0.125}, draw=none, from=3-1, to=5-3]
            \arrow[two heads, from=3-3, to=5-3]
            \arrow["f"', from=5-1, to=5-3]
        \end{tikzcd}
    \end{equation}
    when we vertically paste the two pullback squares and inspect them, we note that $f^\ast W = (f\times \mathsf{id})^\ast W$. The dimension condition is satisfied by inspecting the submersivity conditions.

    \medskip

    Additionally, we must check that the pullback submanifold $f^*W$ satisfies the submersivity conditions~(\ref{id: categorical}) and~(\ref{id: spatial}). This follows because the projection $\pi_{|W}$ is a submersion, and submersions are stable under pullback along smooth maps. 
    
    \medskip

    \nid More explicitly, we have the following commutative diagram:
    \begin{equation}
    \begin{tikzcd}[ampersand replacement=\&]
        \& {(f\times \id)^\ast W_{|}} \&\& {W_{|}} \\
        {(f\times \id)^\ast W} \&\& W \\
        \& {N\times T^{\mathsf{c}}\times V} \&\& {M\times T^{\mathsf{c}}\times V} \\
        {N\times\mathsf{cor}(T)\times V} \&\& {M\times\mathsf{cor}(T)\times V} \\
        \& {N\times T^{\mathsf{c}}} \&\& {M\times T^{\mathsf{c}}}
        \arrow[from=1-2, to=1-4]
        \arrow[from=1-2, to=2-1]
        \arrow[color={rgb,255:red,214;green,153;blue,92}, hook, from=1-2, to=3-2]
        \arrow["\lrcorner"{anchor=center, pos=0.125}, draw=none, from=1-2, to=3-4]
        \arrow["\lrcorner"{anchor=center, pos=0.125, rotate=-90}, draw=none, from=1-2, to=4-1]
        \arrow[from=1-4, to=2-3]
        \arrow[hook, from=1-4, to=3-4]
        \arrow["\lrcorner"{anchor=center, pos=0.125, rotate=-90}, draw=none, from=1-4, to=4-3]
        \arrow[from=2-1, to=2-3]
        \arrow[hook, from=2-1, to=4-1]
        \arrow["\lrcorner"{anchor=center, pos=0.125}, draw=none, from=2-1, to=4-3]
        \arrow[hook, from=2-3, to=4-3]
        \arrow[from=3-2, to=3-4]
        \arrow[from=3-2, to=4-1]
        \arrow[color={rgb,255:red,214;green,153;blue,92}, from=3-2, to=5-2]
        \arrow["\lrcorner"{anchor=center, pos=0.125}, draw=none, from=3-2, to=5-4]
        \arrow[hook', from=3-4, to=4-3]
        \arrow[from=3-4, to=5-4]
        \arrow[from=4-1, to=4-3]
        \arrow["{f\times \id}"', from=5-2, to=5-4]
    \end{tikzcd}
\end{equation}
    This diagram renders manifest the fact that the composite of the two orange arrows is a submersion, since it is the pullback of a submersion along a smooth map, as desired.   
\end{remark}

\subsubsection{Digression on interpretation}

Understanding this definition as a precursor to defining an $(\oo,1)$-category of bordisms requires further unpacking. In \Cref{ch2:: subsec: interval_spaces_model_infinity_cats} we explained that, for interval Segal spaces, the space of $k$-composable morphisms is encoded in the functor's value at a choice of interval $T$ with $\gamma(T)=[k]$. So far, as we explained in the introduction to this section, we do not yet have a pre-interval space, let alone an interval Segal space. Instead, we have a co-presheaf in $\Int$ with values in presheaves on $\Man$ --- taking the transpose under the power-mapping adjunction.

Nevertheless, we can already outline the blueprint for how we would encode a sequence of $k$-composable bordisms
\begin{equation}
    M_1 \xra{\Sigma_1} M_2 \xra{\Sigma_2} \cdots \xra{\Sigma_{k-1}} M_k
\end{equation} 
under this setup. For simplicity and grounding, let us examine the case $d=1$. Choose $T = \sU(A,B,C)\in \Int$ so that $\gamma(T) = [2]$; an element of the set $\widetilde{\bord}_{d,V}(\ast, \sU(A,B,C))$ is a $1$-dimensional submanifold with boundary
\[
W\subset \ast\times \core(T)\times V,
\]
which could be represented as 
\begin{equation}
    \scalebox{0.8}{\tikzfig{figures/Ch3/BOR1}{\begin{tikzpicture}
	\begin{pgfonlayer}{nodelayer}
		\node [style=none] (48) at (2, 11) {};
		\node [style=none] (49) at (5.5, 11) {};
		\node [style=none] (50) at (2, 0) {};
		\node [style=none] (51) at (5.5, 0) {};
		\node [style=none] (0) at (-2, 11) {};
		\node [style=none] (1) at (-2, -0.25) {};
		\node [style=none] (2) at (-2.5, 0) {};
		\node [style=none] (3) at (17.5, 0) {};
		\node [style=none] (4) at (-2.75, 11) {$V$};
		\node [style=none] (5) at (-2, 0) {};
		\node [style=none] (6) at (0, 0) {};
		\node [style=none] (8) at (5.5, 0) {};
		\node [style=none] (9) at (15.5, 0) {};
		\node [style=none] (10) at (17, 0) {};
		\node [style=endpoint] (11) at (17, 0) {};
		\node [style=endpoint] (12) at (-2, 0) {};
		\node [style=none] (13) at (18.5, 0) {$[0,1]$};
		\node [style=none] (14) at (0, 10.25) {};
		\node [style=none] (20) at (0, 2.5) {};
		\node [style=none] (24) at (15.5, 2.25) {};
		\node [style=none] (36) at (1, 10.75) {$W$};
		\node [style=none] (37) at (-2, 11) {};
		\node [style=none] (38) at (0, 11) {};
		\node [style=none] (41) at (15.5, 11) {};
		\node [style=none] (42) at (17, 11) {};
		\node [style=none] (43) at (2, 0) {};
		\node [style=none] (45) at (5.75, 0.5) {};
		\node [style=none] (46) at (2, 11) {};
		\node [style=none] (56) at (15.5, 11) {};
		\node [style=none] (57) at (15.5, 0) {};
		\node [style=none] (58) at (17, 11) {};
		\node [style=none] (59) at (17, 0) {};
		\node [style=none] (53) at (0, 11) {};
		\node [style=none] (55) at (0, 0) {};
		\node [style=none] (52) at (-2, 11) {};
		\node [style=none] (54) at (-2, 0) {};
		\node [style=none] (60) at (8, 11) {};
		\node [style=none] (61) at (12.5, 11) {};
		\node [style=none] (62) at (8, 0) {};
		\node [style=none] (63) at (12.5, 0) {};
		\node [style=none] (68) at (0, 4.25) {};
		\node [style=none] (71) at (15.5, 4) {};
		\node [style=none] (72) at (0, 9.25) {};
		\node [style=none] (73) at (0, 7.75) {};
		\node [style=none] (74) at (2, 9.25) {};
		\node [style=none] (75) at (2, 7.75) {};
		\node [style=none] (76) at (3.25, 8.5) {};
		\node [style=none] (77) at (5.5, 9.25) {};
		\node [style=none] (78) at (5.5, 7.75) {};
		\node [style=none] (79) at (4.25, 8.5) {};
		\node [style=none] (80) at (2, 2.5) {};
		\node [style=none] (81) at (3.75, 3) {};
		\node [style=none] (82) at (2.75, 4.5) {};
		\node [style=none] (83) at (4, 5) {};
		\node [style=none] (84) at (5.5, 5) {};
		\node [style=none] (85) at (8, 5) {};
		\node [style=none] (86) at (10.25, 4.25) {};
		\node [style=none] (87) at (9, 3) {};
		\node [style=none] (88) at (11, 2.75) {};
		\node [style=none] (89) at (12.5, 2.75) {};
		\node [style=none] (90) at (8, 9.25) {};
		\node [style=none] (91) at (8, 7.75) {};
		\node [style=none] (92) at (9.5, 8.5) {};
	\end{pgfonlayer}
	\begin{pgfonlayer}{edgelayer}
		\draw [style=red shade low op] (51.center)
			 to (50.center)
			 to (48.center)
			 to (49.center);
		\draw [style=dir] (1.center) to (0.center);
		\draw [style={open_int}] (5.center) to (6.center);
		\draw [style={open_int}] (9.center) to (10.center);
		\draw [style=red shade low op] (41.center) to (9.center);
		\draw [style=red shade low op] (9.center) to (11);
		\draw [style=red shade low op] (11) to (42.center);
		\draw [style=red shade low op] (42.center) to (41.center);
		\draw [style={open_int}] (43.center) to (8.center);
		\draw [style=red shade low op] (59.center)
			 to (58.center)
			 to (56.center)
			 to (57.center)
			 to cycle;
		\draw [style=dir] (2.center) to (3.center);
		\draw [style=red shade low op] (55.center)
			 to (54.center)
			 to (52.center)
			 to (53.center);
		\draw [style={open_int}] (43.center) to (8.center);
		\draw [style={open_int}] (54.center) to (55.center);
		\draw [style={open_int}] (57.center) to (59.center);
		\draw [style=dir] (54.center) to (52.center);
		\draw [style=red shade low op] (63.center)
			 to (62.center)
			 to (60.center)
			 to (61.center);
		\draw [style=geometric, in=180, out=0] (72.center) to (74.center);
		\draw [style=geometric, in=180, out=15] (73.center) to (75.center);
		\draw [style=geometric, in=90, out=0] (74.center) to (76.center);
		\draw [style=geometric, in=270, out=0] (75.center) to (76.center);
		\draw [style=geometric, in=90, out=180] (77.center) to (79.center);
		\draw [style=geometric, in=-90, out=180] (78.center) to (79.center);
		\draw [style=geometric] (20.center) to (80.center);
		\draw [style=geometric, in=-90, out=0, looseness=0.75] (80.center) to (81.center);
		\draw [style=geometric, in=-90, out=90] (81.center) to (82.center);
		\draw [style=geometric, in=-180, out=90, looseness=0.75] (82.center) to (83.center);
		\draw [style=geometric] (83.center) to (84.center);
		\draw [style=geometric] (84.center) to (85.center);
		\draw [style=geometric, in=105, out=0, looseness=0.75] (85.center) to (86.center);
		\draw [style=geometric, in=90, out=-90, looseness=0.75] (86.center) to (87.center);
		\draw [style=geometric, in=-180, out=-90, looseness=0.50] (87.center) to (88.center);
		\draw [style=geometric, in=180, out=0] (88.center) to (89.center);
		\draw [style=geometric, in=-180, out=0] (89.center) to (71.center);
		\draw [style=geometric] (78.center) to (91.center);
		\draw [style=geometric] (77.center) to (90.center);
		\draw [style=geometric, in=90, out=0] (90.center) to (92.center);
		\draw [style=geometric, in=270, out=0] (91.center) to (92.center);
		\draw [style={open_int}] (62.center) to (63.center);
	\end{pgfonlayer}
\end{tikzpicture}
}}
\end{equation}
\smallskip

\noindent In which case we have a $1$-dimensional submanifold $W\subset \ast \times \sU(A,B,C) \times V$ where 
\begin{equation}
\sU(A,B,C) = \scalebox{0.6}{\tikzfig{figures/Ch3/BOR2}{\begin{tikzpicture}
	\begin{pgfonlayer}{nodelayer}
		\node [style=none] (50) at (-1.75, 0) {};
		\node [style=none] (51) at (1.75, 0) {};
		\node [style=none] (5) at (-5.75, 0) {};
		\node [style=none] (6) at (-3.75, 0) {};
		\node [style=none] (8) at (1.75, 0) {};
		\node [style=none] (9) at (5.25, 0) {};
		\node [style=none] (10) at (6.75, 0) {};
		\node [style=endpoint] (11) at (6.75, 0) {};
		\node [style=endpoint] (12) at (-5.75, 0) {};
		\node [style=none] (43) at (-1.75, 0) {};
		\node [style=none] (57) at (5.25, 0) {};
		\node [style=none] (59) at (6.75, 0) {};
		\node [style=none] (55) at (-3.75, 0) {};
		\node [style=none] (54) at (-5.75, 0) {};
	\end{pgfonlayer}
	\begin{pgfonlayer}{edgelayer}
		\draw (54.center) to (59.center);
		\draw [style=red shade] (51.center) to (50.center);
		\draw [style={open_int}] (5.center) to (6.center);
		\draw [style={open_int}] (9.center) to (10.center);
		\draw [style=red shade] (9.center) to (11);
		\draw [style={open_int}] (43.center) to (8.center);
		\draw [style=red shade] (57.center) to (59.center);
		\draw [style=red shade] (55.center) to (54.center);
		\draw [style={open_int}] (43.center) to (8.center);
		\draw [style={open_int}] (54.center) to (55.center);
		\draw [style={open_int}] (57.center) to (59.center);
	\end{pgfonlayer}
\end{tikzpicture}
}}.
\end{equation}
The only non-trivial submersivity condition is (\ref{id: categorical}), and it translates to the projection
\begin{equation}
    \pi \maps \big(W\cap (\ast \times (A\cup B \cup C)\times V)\big) \fib \ast \times (A\cup B \cup C)
\end{equation}
onto the closed intervals $A$, $B$, and $C$ being a submersion. Recall from the boundary version of Ehresmann's fibration theorem, \Cref{ap:: prop: ehresmann boundary}, that proper submersions are smooth fiber bundles. Moreover, in this case the targets are closed intervals which are contractible, so from \Cref{ap:: prop: triviality contractibility} we have that 
\begin{equation}
    W_{|A_1\sqcup A_2 \sqcup A_3} \simeq (W_1 \sqcup W_2 \sqcup W_3) \times [0,1]. 
\end{equation}    
On the other hand, over the intervals $A_i$ (depicted in red) there are no submersivity conditions so the projection restricted to these open subsets can admit any assortment of critical points. 

These red regions, where the submanifold is allowed to change its topology, store the content of morphisms. In this chapter, these data are $d$-dimensional unoriented bordisms. Later in this chapter and, more importantly, in Chapter~5, we will generalize this construction to include additional structure, i.e.\ fiberwise tangential structures and stratifications. 

A few more words are still deserved. We note that the non-shaded regions are effectively storing data for the sources and targets of our morphisms, already with a ``baked-in'' choice of collar.  

We end this digression with a discussion on spatial structure. More concretely, we address how and where this construction encodes the idea that higher morphisms, or invertible morphisms in the $(\oo,1)$-category $\bord_d$, should be diffeomorphisms of $d$-dimensional bordisms, and homotopies should be isotopies between such diffeomorphisms, etc. Paths will correspond to concordances of the point $\ast$, more concretely to sections for $M=\AA^1$. A path is thus a $(d+1)$-dimensional smooth submanifold $\Sigma \subset \AA^1 \times \sU(A,B,C)\times V$, for example 
\begin{equation}
    \tikzfig{figures/Ch3/BOR12}{\begin{tikzpicture}
	\begin{pgfonlayer}{nodelayer}
		\node [style=endpoint] (23) at (9, 3) {};
		\node [style=none] (0) at (-3, 3) {};
		\node [style=none] (1) at (9.5, 3) {};
		\node [style=none] (2) at (-1, 3) {};
		\node [style=none] (3) at (8, 3) {};
		\node [style=none] (5) at (-0.5, 3) {};
		\node [style=none] (6) at (0.75, 3) {};
		\node [style=none] (7) at (5.25, 3) {};
		\node [style=none] (8) at (6.75, 3) {};
		\node [style=endpoint] (12) at (-3, 3) {};
		\node [style=none] (18) at (1.775, 3) {};
		\node [style=none] (20) at (9, 3) {};
		\node [style=none] (21) at (4, -5) {};
		\node [style=none] (22) at (9, 3) {};
		\node [style=endpoint] (24) at (4, -5) {};
		\node [style=none] (25) at (-7.95, -5) {};
		\node [style=none] (26) at (4.5, -5) {};
		\node [style=none] (27) at (-6, -5) {};
		\node [style=none] (28) at (3, -5) {};
		\node [style=none] (29) at (-5.5, -5) {};
		\node [style=none] (30) at (-4.25, -5) {};
		\node [style=none] (31) at (1.25, -5) {};
		\node [style=none] (32) at (1.75, -5) {};
		\node [style=none] (33) at (-5.75, -5.75) {$A$};
		\node [style=none] (34) at (-2.125, -5.75) {$B$};
		\node [style=none] (35) at (2, -5.75) {$C$};
		\node [style=none] (45) at (4, -5) {};
		\node [style=none] (46) at (-3, 10) {};
		\node [style=none] (47) at (-3, 3) {};
		\node [style=none] (48) at (-8.25, -5.5) {};
		\node [style=none] (49) at (-3, 3) {};
		\node [style=none] (71) at (3.75, -5.5) {};
		\node [style=none] (72) at (9, 3) {};
		\node [style=none] (73) at (4.5, 3) {};
		\node [style=none] (75) at (-0.25, 3) {};
		\node [style=none] (76) at (6.25, 3) {};
		\node [style=none] (77) at (1.25, -5) {};
		\node [style=none] (86) at (-4.5, 10) {$\AA^\bullet \times V$};
		\node [style=none] (87) at (-9, -5.5) {$\AA^1$};
		\node [style=none] (88) at (0.75, 3) {};
		\node [style=none] (89) at (2, 3) {};
		\node [style=none] (90) at (-0.5, -5) {};
		\node [style=none] (91) at (1.25, -5) {};
		\node [style=none] (92) at (2, 6.3) {};
		\node [style=none] (50) at (0.75, 6) {};
		\node [style=none] (51) at (5.25, 8) {};
		\node [style=none] (69) at (-4.25, -2) {};
		\node [style=none] (70) at (1.25, -4.5) {};
	\end{pgfonlayer}
	\begin{pgfonlayer}{edgelayer}
		\draw [style=ts-dashed] (50.center) to (6.center);
		\draw [style=dir] (0.center) to (1.center);
		\draw [style={open_int}] (3.center) to (22.center);
		\draw [style={open_int}] (0.center) to (2.center);
		\draw [style={open_int}] (5.center) to (6.center);
		\draw [style={open_int}] (7.center) to (8.center);
		\draw [style=blue shade] (21.center) to (20.center);
		\draw [style=blue shade] (20.center) to (18.center);
		\draw [style=red shade] (3.center) to (22.center);
		\draw [style=red shade] (3.center) to (22.center);
		\draw [style=red shade] (8.center) to (3.center);
		\draw [style=dir] (25.center) to (26.center);
		\draw [style={open_int}] (28.center) to (45.center);
		\draw [style={open_int}] (25.center) to (27.center);
		\draw [style={open_int}] (29.center) to (30.center);
		\draw [style={open_int}] (31.center) to (32.center);
		\draw [style=red shade] (28.center) to (45.center);
		\draw [style=red shade] (28.center) to (45.center);
		\draw [style=red shade] (32.center) to (28.center);
		\draw [style=dir] (47.center) to (48.center);
		\draw [style=dir] (49.center) to (46.center);
		\draw [style=red shade] (77.center) to (76.center);
		\draw [style=red shade] (76.center) to (49.center);
		\draw [style=red shade, in=135, out=45, loop] (49.center) to ();
		\draw [style=red shade] (49.center) to (73.center);
		\draw [style=blue shade] (45.center) to (72.center);
		\draw [style=blue shade] (72.center) to (75.center);
		\draw [style=red shade] (76.center) to (77.center);
		\draw [style=red shade] (77.center) to (25.center);
		\draw [style=red shade] (25.center) to (49.center);
		\draw [style=red shade] (49.center) to (76.center);
		\draw [style=ts-dashed] (69.center) to (30.center);
		\draw [style=ts-dashed] (51.center) to (7.center);
		\draw [style=ts-dashed, in=90, out=-90] (70.center) to (31.center);
		\draw [style=geometric shaded] (69.center)
			 to [in=165, out=0] (70.center)
			 to (51.center)
			 to [in=0, out=-180] (50.center)
			 to [in=60, out=-120] cycle;
	\end{pgfonlayer}
\end{tikzpicture}
}.
\end{equation}
Condition (\ref{id: categorical}) ensures that the collars behave as expected, and condition (\ref{id: spatial}) enforces the idea that paths of composable sequences of bordisms are diffeomorphisms of the composed bordism which intertwine with the decomposition. This can be seen explicitly by inspecting the fibers over the inner components.

The next few results establish the locality properties of $\widetilde{\bord}_{d,V}$ needed for the construction of our main object. In turn, these are what allow the resulting interval Segal space to have the correct symmetric monoidal homotopy category. 

\begin{proposition}\label{ch3:: prop: sheaf condition}
    The adjoint functor 
    \begin{equation}
    \widetilde{\bord}_{d,V}(-,-)\maps \Int \longra \Fun(\man^{\op},\sset)
    \end{equation}
    factors through $\shv(\man)$, the category of $\oo$-sheaves on $\man$ relative to the Grothendieck topology of open covers. In other words, for every $T\in \Int$ the functor
    \[
    \widetilde{\bord}_{d,V}(-,T)\maps \man^{\op}\longra \set
    \]
    is a smooth $\oo$-sheaf of (discrete) simplicial sets.
\end{proposition}

\begin{proof}
    We will check conditions (\ref{binary cover}) and (\ref{disj union}) from \Cref{ch3:: thm: BBEP-sheaf-criterion}. We begin with (\ref{binary cover}): let $M$ be a smooth manifold and let $\{U_1,U_2\}$ be an open cover of $M$. We must show that the \emph{decomposition} map
    \begin{equation}
        \label{eq1}
        \widetilde{\bord}_{d,V}(M,T) \longrightarrow \widetilde{\bord}_{d,V}(U_1,T)\underset{\widetilde{\bord}_{d,V}(U_1\cap \,U_2,T)}{\times}\widetilde{\bord}_{d,V}(U_2,T)
    \end{equation}
    is an isomorphism. Injectivity is almost immediate: given a pair of different submanifolds with boundary 
    \[
    W,W' \subset \big(M\times \core(T)\times V\big),
    \]
    their images under the decomposition map \eqref{eq1} must differ, manifestly, in at least one of the factors of the pullback.
    
     The surjectivity of the decomposition map is downstream of the fact that our construction works with submanifolds embedded into some larger universe, as opposed to a set (or even space) of manifolds --- in this case $M \times \core(T) \times V$. An element of the codomain of \eqref{eq1} is thus equivalent to the datum of
    \begin{equation}
    \big(W_1 \subset U_1 \times \mathsf{cor}(T)\times V\big), \quad \big(W_2 \subset U_2 \times \mathsf{cor}(T)\times V \big),
    \end{equation} 
    a pair of submanifolds with boundary satisfying conditions (\ref{id: categorical}) and (\ref{id: spatial}) of \Cref{ch3:: def: bord_d}, whose restrictions along $\big(U_1 \cap U_2\big)\times \core(T)\times V$ agree pointwise --- both equal to the submanifold
    \begin{equation}
    W_{|U_{1}\cap U_{2}}\subset \big( U_1\cap U_2\big) \times \mathsf{cor}(T)\times V.
    \end{equation}
    To construct an element in $\widetilde{\bord}_{d,V}(M,T)$ which decomposes to this data, we glue, i.e.\ take the following pushout
    \begin{equation}
        \begin{tikzcd}
            {W_{|U_{1}\cap\,U_{2}}} && {W_{2}} \\
            \\
            {W_{1}} && {W_{1}\,\underset{W_{|U_{1}\cap\,U_{2}}}{\cup} \, W_2}
            \arrow[hook, from=1-1, to=1-3]
            \arrow[hook, from=1-1, to=3-1]
            \arrow[hook, from=1-3, to=3-3]
            \arrow[hook, from=3-1, to=3-3]
            \arrow["\lrcorner"{anchor=center, pos=0.125, rotate=180}, draw=none, from=3-3, to=1-1]
        \end{tikzcd}
    \end{equation}
    in $\man$. This pushout restricts to $W_1$ and $W_2$ on the open subsets $U_1$ and $U_2$, respectively, and satisfies conditions (\ref{id: categorical}) and (\ref{id: spatial}). The latter follows from the fact that submersion is an open condition, and the pushout was carried out along an open cover. 

    Condition (\ref{disj union}) from \Cref{ch3:: thm: BBEP-sheaf-criterion} is checked similarly, using the fact that disjoint unions of manifolds correspond to coproducts in $\man$ and that the construction is taken componentwise.
\end{proof}
\begin{remark}
    The simplicial set $\bord_{d,V}(\AA^{\bullet},T)$ is a $0$-weak Kan complex for every $T\in \Int$. This follows from the general fact that for every concordance-invariant sheaf $\cF$, the simplicial set 
    \begin{equation}
        \big | \cF(\AA^{\bullet} \times *)\big | \in \sset
    \end{equation}
    is a $0$-weak Kan complex.
\end{remark}

Having entertained the locality of our construction with respect to the smooth manifold variable, we now assert yet another form of locality; again fixing $T \in \Int$, the next proposition proves a ``naive sheaf'' property with respect to binary open covers of the closed interval $[0,1]$ itself.

\begin{proposition}[Naive sheaf condition in the interval direction]\label{ch3:: prop: naive sheaf con}
Let $T\in \Int$, and let $\{U,V\}$ be an open cover of $\core(T)\subset [0,1]$. For each $k\geq 0$, the canonical map
\begin{equation}
\label{eq:naive sheaf interval levelwise}
\bord_d\big(\AA^{k},T\big)\longrightarrow
\bord_d\big(\AA^{k},T\big)_{|U}\underset{\bord_d\big(\AA^{k},T\big)_{|U\cap V}}{\times}
\bord_d\big(\AA^{k},T\big)_{|V}
\end{equation}
is a bijection of sets. Equivalently, the simplicial set $\bord_d(\AA^{\bullet},T)$ is the strict pullback
\begin{equation}
\label{eq:naive sheaf interval simplicial}
\bord_d(\AA^{\bullet},T)\;\cong\;
\bord_d(\AA^{\bullet},T)_{|U}\underset{\bord_d(\AA^{\bullet},T)_{|U\cap V}}{\times}
\bord_d(\AA^{\bullet},T)_{|V},
\end{equation}
where
\begin{equation}
\label{eq:bord-restriction-U}
\bord_d(\AA^{k},T)_{|U}
=
\Big\{
W\subseteq \AA^{k}\times\big(U\cap \core(T)\big)\times V
\Big\}.
\end{equation}
Here $\bord_d(\AA^{k},T)_{|U}$ denotes the set of all $(d+k)$-dimensional
submanifolds with boundary
\[
W\subseteq \AA^{k}\times\big(U\cap \core(T)\big)\times V
\]
which satisfy conditions (\ref{id: dim}),(\ref{id: categorical}), and (\ref{id: spatial}) from \Cref{ch3:: def: bord_d}, and the restriction maps are given by intersecting with
\[
\AA^{k}\times\big(U\cap V\cap \core(T)\big)\times V\subseteq 
\AA^{k}\times\big(U\cap\core(T)\big)\times V,
\]
and similarly for $V$.
\end{proposition}

\begin{proof}
Fix $k\geq 0$ and abbreviate
\begin{equation}
Z = \AA^{k}\times \core(T)\times V.
\end{equation}
An element of $\bord_d(\AA^{k},T)$ is a $(d+k)$–dimensional submanifold with boundary
\begin{equation}
W\subseteq Z
\end{equation}
satisfying conditions (\ref{id: dim})--(\ref{id: spatial}) of \Cref{ch3:: def: bord_d}. Its images in
$\bord_d(\AA^{k},T)_{|U}$ and $\bord_d(\AA^{k},T)_{|V}$ are given by intersection with
$\AA^{k}\times\big(U\cap \core(T)\big)\times V$ and
$\AA^{k}\times\big(V\cap \core(T)\big)\times V$, respectively, so injectivity of
\eqref{eq:naive sheaf interval levelwise} is immediate.

For surjectivity, start with a pair of submanifolds
\begin{equation}
W_U\subseteq \AA^{k}\times\big(U\cap\core(T)\big)\times V,\qquad
W_V\subseteq \AA^{k}\times\big(V\cap\core(T)\big)\times V,
\end{equation}
each satisfying (\ref{id: dim})--(\ref{id: spatial}), whose restrictions to
\begin{equation}
\AA^{k}\times\big(U\cap V\cap\core(T)\big)\times V
\end{equation}
coincide as subsets. Define
\begin{equation}
W := W_U \cup W_V \subseteq Z.
\end{equation}
Since $U$ and $V$ cover $\core(T)$ and $W_U$, $W_V$ agree on the overlap, $W$ is a well-defined subset of $Z$.

Standard differential-topological arguments ensure that $W$ is a smooth
submanifold with boundary: all the structure is determined locally, and at each point of $W$ either $W$ coincides with $W_U$, or with $W_V$, or with their common restriction in the overlap. Condition (\ref{id: dim}) is local, so it also holds for $W$.

The submersivity conditions (\ref{id: categorical}) and (\ref{id: spatial}) are likewise local on the source: they hold on $W_U$ and $W_V$ by assumption, hence they hold on $W$, since submersion is an open condition. This shows that $W$ lies in $\bord_d(\AA^{k},T)$ and restricts to $(W_U,W_V)$, proving surjectivity and hence the bijectivity of \eqref{eq:naive sheaf interval levelwise}. The statement for the simplicial set follows levelwise.
\end{proof}

We now come to the main idea of this section. We will use the concordance resolution of sheaves, restricted to the smooth cosimplicial object, in order to construct an interval Segal space. For each interval object $I$, the functor
\begin{equation}
    \widetilde{\bord}_{d,V}(-,I) \maps \man^{\op}\longra \mathsf{Set}
\end{equation}
is set-valued, so evaluating it on the smooth cosimplicial simplex already produces a simplicial set
\begin{equation}
    \widetilde{\bord}_{d,V}(\AA^\bullet,I).
\end{equation}
Thus, in the present case, the concordance-resolution construction introduces no additional simplicial direction beyond this one. The next result is both a definition and a theorem.

\medskip

\begin{theorem}\label{ch3:: thm: interval segal bordism cat}
    Let $d \geq 0$ be a non-negative integer and let $V$ be a finite-dimensional vector space of sufficiently high dimension. The composite functor
    \begin{align}
        \bord_{d,V} \maps \Int &\longra \shv(\man)  \longra \sset  \\
         I &\longmapsto \widetilde{\bord}_{d,V}(-,I) \longmapsto \bB\widetilde{\bord}_{d,V}(*,I) = \big|\,\widetilde{\bord}_{d,V}(\AA^\bullet,I)\,\big| 
    \end{align}
    is an interval Segal space.
\end{theorem}

\smallskip

\begin{proof}
    \textbf{Homotopical property.}
    We begin by checking that the functor sends $\pi_0$-equivalences to weak equivalences in $\sset$. Concretely, for every open embedding $T\hookra S$ such that $\pi_0(T) \simeq \pi_0(S)$, the associated map of simplicial sets 
    \begin{equation}
        \bord_{d,V}(T) \xra{\,\,\,\,\mathsf{forget}_{T \ra S}\,\,\,\,} \bord_{d,V}(S),
    \end{equation}   
    is a weak equivalence. There is a map
    \[
    \mathsf{ext}\maps \bord_{d,V}(S) \longra \bord_{d,V}(T)
    \]
    which ``squeezes'' the underlying submanifold with boundary, in particular bunching up critical points. Let 
    \begin{equation}
        T = U_1 \cup \cdots \cup U_k, \quad S = \widetilde{U}_1 \cup \cdots \cup \widetilde{U}_k,
    \end{equation} 
    with $0\in U_1,\widetilde{U}_1$ and $1\in U_k,\widetilde{U}_k$ --- note that $U_i \subseteq \widetilde{U}_i$. Let $s\maps [0,1]\longra [0,1]$ be a diffeomorphism which maps $\widetilde{U}_i \mapsto U_i$ and is the identity elsewhere; essentially a rescaling map obtained by smoothing a piecewise-linear map. We now define
    \begin{align*}
        \bord_{d,V}(S) &\xra{\,\,\,\mathsf{ext}_{S \longra T}\,\,\,} \bord_{d,V}(T) \\
        \big( W \subset \AA^\bullet \times [0,1]\times V \big) &\longmapsto \big((\id_{\AA^\bullet} \times\, s \times \id_{V})^{*}W\subset \AA^\bullet \times [0,1]\times V \big).
    \end{align*}
    The map is manifestly well-defined since $(\id_{\AA^\bullet} \times\, s \times \id_{V})$ is a diffeomorphism, i.e.\ it preserves condition (\ref{id: dim}) and the submersivity conditions (\ref{id: categorical}) and (\ref{id: spatial}) over $T$. 

    Neither of the two composites is strictly the identity, so we must articulate appropriate simplicial homotopies witnessing the desired equivalence. To do so, we observe that a simplicial homotopy 
    \begin{equation}
        \bord_{d,V}(S)\times \Delta^1 \longra \bord_{d,V}(S)
    \end{equation} 
    amounts, via adjunction, to a simplicial map 
    \begin{equation}
        \widetilde{\bord}_{d,V}(\AA^\bullet,S) \longra \widetilde{\bord}_{d,V}(\AA^\bullet \times \AA^1, S)
    \end{equation}
    whose further restriction to $\{0\}$ and $\{1\}$ yields maps identical to $\big(\mathsf{forget}\circ\mathsf{ext}\big)$ and $\id$, respectively. We construct this map by pulling back along 
    \begin{align*}
        \big(\AA^\bullet \times (\AA^1 \times [0,1]) \times V\big) \xra{\id_{\AA^\bullet}\times h \times \id_{V}} \big( \AA^\bullet \times  [0,1]\times V \big),
    \end{align*}
    where $h\maps \AA^1 \times [0,1] \longra [0,1]$, defined as $h(t,x) = (1-t) \cdot s(x) + t \cdot x$, is a smooth linear isotopy between the diffeomorphism $s$ and the identity --- whose graph we depict in the figure below: 
    \begin{equation}
        \scalebox{0.8}{\tikzfig{figures/Ch3/BOR14}{\begin{tikzpicture}
	\begin{pgfonlayer}{nodelayer}
		\node [style=none] (29) at (-6, -6.25) {};
		\node [style=none] (30) at (6.5, -6.25) {};
		\node [style=none] (31) at (-6, 6.25) {};
		\node [style=none] (35) at (6, -6.25) {};
		\node [style=none] (55) at (-7.25, 6.5) {$[0,1]$};
		\node [style=none] (57) at (7.5, -6.5) {$[0,1]$};
		\node [style=none] (58) at (-3.5, -6.25) {};
		\node [style=none] (59) at (-6, -5.25) {};
		\node [style=none] (60) at (-1, -6.25) {};
		\node [style=none] (61) at (1, -6.25) {};
		\node [style=none] (62) at (2, -6.25) {};
		\node [style=none] (63) at (3.5, -6.25) {};
		\node [style=none] (64) at (-6, 5.75) {};
		\node [style=none] (65) at (-3.5, -5.25) {};
		\node [style=none] (66) at (-6, -0.75) {};
		\node [style=none] (67) at (-6, 0.25) {};
		\node [style=none] (68) at (-6, 2.75) {};
		\node [style=none] (69) at (-6, 3.75) {};
		\node [style=none] (70) at (4.5, -6.25) {};
		\node [style=none] (71) at (-6, 4.75) {};
		\node [style=none] (72) at (-1, -0.75) {};
		\node [style=none] (73) at (1, 0.25) {};
		\node [style=none] (74) at (2, 2.75) {};
		\node [style=none] (75) at (3.5, 3.75) {};
		\node [style=none] (76) at (4.5, 4.75) {};
		\node [style=none] (77) at (6, 5.75) {};
		\node [style=none] (81) at (-4.25, -5.5) {};
	\end{pgfonlayer}
	\begin{pgfonlayer}{edgelayer}
		\draw [style=dir] (29.center) to (30.center);
		\draw [style=dir] (29.center) to (31.center);
		\draw [style={open_int}] (29.center) to (58.center);
		\draw [style={open_int_blue}] (29.center) to (59.center);
		\draw (29.center) to (65.center);
		\draw [style={open_int}] (60.center) to (61.center);
		\draw [style={open_int_blue}] (66.center) to (67.center);
		\draw [style={open_int}] (62.center) to (63.center);
		\draw [style={open_int_blue}] (68.center) to (69.center);
		\draw [style={open_int}] (70.center) to (35.center);
		\draw [style={open_int_blue}] (64.center) to (71.center);
		\draw (72.center) to (73.center);
		\draw (74.center) to (75.center);
		\draw (76.center) to (77.center);
		\draw [style=new edge style 0, in=-150, out=30] (65.center) to (72.center);
		\draw [style=new edge style 0, in=-135, out=30, looseness=0.50] (73.center) to (74.center);
		\draw [style=new edge style 0, in=-150, out=30, looseness=0.75] (75.center) to (76.center);
		\draw [style=geometric] (29.center) to (77.center);
		\draw [style=ts-dashed-grey] (59.center) to (65.center);
		\draw [style=ts-dashed-grey] (65.center) to (58.center);
		\draw [style=ts-dashed-grey] (66.center) to (72.center);
		\draw [style=ts-dashed-grey] (72.center) to (60.center);
		\draw [style=ts-dashed-grey] (61.center) to (73.center);
		\draw [style=ts-dashed-grey] (73.center) to (67.center);
		\draw [style=ts-dashed-grey] (68.center) to (74.center);
		\draw [style=ts-dashed-grey] (74.center) to (62.center);
		\draw [style=ts-dashed-grey] (63.center) to (75.center);
		\draw [style=ts-dashed-grey] (75.center) to (69.center);
		\draw [style=ts-dashed-grey] (71.center) to (76.center);
		\draw [style=ts-dashed-grey] (76.center) to (70.center);
		\draw [style=ts-dashed-grey] (64.center) to (77.center);
		\draw [style=ts-dashed-grey] (77.center) to (35.center);
	\end{pgfonlayer}
\end{tikzpicture}
}}
    \end{equation}
    Conditions (\ref{id: dim}), (\ref{id: categorical}), and (\ref{id: spatial}) are preserved since the isotopy is a diffeotopy, which ensures a well-defined simplicial map. 
    
    To establish that $(\mathsf{ext}\circ \mathsf{forget}) \simeq \id_{\bord_{d,V}(T)}$ we construct the analogous simplicial homotopy 
    \begin{equation}
    \widetilde{\bord}_{d,V}(T)\times \Delta^1 \longra \widetilde{\bord}_{d,V}(T)
    \end{equation}
    this time using $s^{-1}(x)$ in place of $s(x)$ when describing the homotopy $h$.

    \medskip

    \textbf{Interval Segal condition.}
    Recall that the interval Segal condition (\Cref{ch2:: def: interval Segal condition}) states that the following diagram of simplicial sets\footnote{These happen to be $0$-weak Kan complexes, and not entirely arbitrary simplicial sets.} is a homotopy pullback square for every $\sU(A,B,C)\in \cJ_2$:
    \begin{equation}\label{eq:Segal condition for bordism}
        \begin{tikzcd}[ampersand replacement=\&]
            {\bord_d\big(U(A,B,C)\big)} \&\& {\bord_d\big(U(A,B)\big)} \\
            \\
            {\bord_d\big(U(B,C)\big)} \&\& {\bord_d\big(U(B)\big)}
            \arrow[from=1-1, to=1-3]
            \arrow[from=1-1, to=3-1]
            \arrow["h"{description, pos=0.3}, "\lrcorner"{anchor=center, pos=0.125}, draw=none, from=1-1, to=3-3]
            \arrow[from=1-3, to=3-3]
            \arrow[from=3-1, to=3-3]
        \end{tikzcd}
    \end{equation}
     Proving that the diagram above is a homotopy pullback square is equivalent to showing that the following diagram is a homotopy pullback square in $\sset$: 
    \begin{equation}\label{eq pullback 1}
        \begin{tikzcd}[ampersand replacement=\&]
            {\bord_d\big(U(A,B,C)\big)} \&\& {\bord_d\big(U(A,B)\big)\times \bord_d\big(U(B,C)\big)} \\
            \\
            {\Map_{\sset}\Big(\Delta^1,\bord_d\big(U(B)\big)\Big)} \&\& {\Map_{\sset}\Big(\partial\Delta^1,\bord_d\big(U(B)\big)\Big)}
            \arrow[from=1-1, to=1-3]
            \arrow[from=1-1, to=3-1]
            \arrow["h"{description, pos=0.3}, "\lrcorner"{anchor=center, pos=0.125}, draw=none, from=1-1, to=3-3]
            \arrow[from=1-3, to=3-3]
            \arrow[from=3-1, to=3-3]
        \end{tikzcd}
    \end{equation}
    Note that all we did was apply a Kan fibration replacement. We now take advantage of the geometric nature of our construction to re-express the homotopy-theoretic path fibration using concordances: namely, we use the fact that, for all $T\in \Int$, the $\oo$-sheaf 
    \begin{equation}
    |\widetilde{\bord}_{d,V}(-\times \AA^{\bullet},T)| \maps \man^{\op} \longra \sset 
    \end{equation}
    is concordance-invariant, i.e.\ for every $M$ the projection
    \[
    |\widetilde{\bord}_{d,V}(M\times \AA^{\bullet},T)| \xra{\,\,\simeq\,\,} |\widetilde{\bord}_{d,V}(M\times \AA^{\bullet}\times \AA^1,T)|
    \]
    induces a weak equivalence of simplicial sets. Thereby, paths in $\bord_d(T)$ are equivalent to concordances in $\bord_d(T)$. We can extend the original diagram~(\ref{eq pullback 1}) by a homotopy equivalent one as follows:
    \begin{equation}
        \begin{tikzcd}[ampersand replacement=\&]
            {\bord_d\big(U(A,B,C)\big)} \&\& {\bord_d\big(U(A,B)\big)\times \bord_d\big(U(B,C)\big)} \\
            \\
            {\mathsf{Map}_{\sset}\Big(\Delta^1,\bord_d\big(U(B)\big)\Big)} \&\& {\mathsf{Map}_{\sset}\Big(\partial\Delta^1,\bord_d\big(U(B)\big)\Big)} \\
            \\
            {|\widetilde{\bord}_{d,V}\big(\AA^{\bullet}\times \AA^1,U(B)\big)|} \&\& {\bord_d\big(U(B)\big)\times \bord_d\big(U(B)\big)}
            \arrow[from=1-1, to=1-3]
            \arrow[from=1-1, to=3-1]
            \arrow["\lrcorner"{anchor=center, pos=0.125, rotate=45}, draw=none, from=1-1, to=3-3]
            \arrow[from=1-3, to=3-3]
            \arrow[from=3-1, to=3-3]
            \arrow["\simeq"', from=3-1, to=5-1]
            \arrow[equals, from=3-3, to=5-3]
            \arrow["{\text {weak Kan fibration}}"', from=5-1, to=5-3]
        \end{tikzcd}
    \end{equation} 
    where the left bottom vertical arrow is a weak equivalence precisely due to the concordance-invariance of $\bord_d\left(\sU(B)\right)$. This means that the homotopy limit of the outer span is weakly equivalent to the homotopy limit of the inner span
    \begin{equation}
        \begin{tikzcd}[ampersand replacement=\&]
            {\mathsf{holim}\Big[\bord_d\big(U(A,B)\big)\times \bord_d\big(U(B,C)\big) \ra \bord_d (U(B))^{\times 2} \la \bord_d\big(\AA^{\bullet}\times \AA^1,U(B)\big)\Big]} \\
            \\
            {\mathsf{holim}\Big[\bord_d\big(U(A,B)\big)\times \bord_d\big(U(B,C)\big) \ra \bord_d (U(B))^{\times 2}\la \mathsf{Map}_{\sset}\Big(\Delta^1, \bord_d\big(U(B)\big)\Big)\Big]}.
            \arrow[from=1-1, to=3-1]
        \end{tikzcd}
    \end{equation}
    Therefore, the interval Segal condition (\ref{eq:Segal condition for bordism}) holds if and only if the following map is a weak equivalence:
    \begin{equation}\label{ch3:: eq: segal condition 1}
        \bord_{d}\big(U(A,B,C)\big) \xra{\,\, \simeq \,\,} \Big(\bord_d\big(U(A,B)\big)\times \bord_d\big(U(B,C)\big)\Big) \underset{\bord_d\big(U(B)\big)\times \bord_d\big(U(B)\big)}{\times^h} \bord_d\big(\AA^{\bullet}\times \AA^1,U(B)\big).
    \end{equation}


    Using the fact (proved in \Cref{ch3:: prop: weak-kan-endpoints}) that for concordance-invariant $\oo$-sheaves, the endpoint map 
    \begin{equation}
    |\widetilde{\bord}_{d,V}(\AA^\bullet\times \AA^1,U(B))| \longra |\widetilde{\bord}_{d,V}(\AA^\bullet \times \partial \AA^1, U(B))|
    \end{equation} 
    is a weak Kan fibration, together with the fact that weak Kan fibrations are realization fibrations (\Cref{ch3:: thm: weak-kan-realization}), we see that showing the interval Segal condition is equivalent to showing that the following map is a weak equivalence of simplicial spaces with the injective model structure:
    \begin{equation}
    \widetilde{\bord}_{d,V}\big(\AA^{\bullet}, U(A,B,C)\big) \xra{\,\, \simeq \,\,} \Big(\widetilde{\bord}_{d,V}\big(\AA^{\bullet},U(A,B)\big)\times \widetilde{\bord}_{d,V}\big(\AA^{\bullet},U(B,C)\big)\Big) \underset{\widetilde{\bord}_{d,V}\big(\AA^{\bullet},U(B)\big)^{\times 2}}{\times^h} \widetilde{\bord}_{d,V}\big(\AA^{\bullet}\times \AA^1,U(B)\big).
    \end{equation}  

    \noindent Alternatively, the right-hand side of the above diagram can be expressed as the following iterated (strict) fiber product:
    \begin{equation}
        \label{eq:iterated fiber product for Segal condition 1}
        \widetilde{\bord}_{d,V}\big(\AA^{\bullet},U(A,B)\big) \underset{\widetilde{\bord}_{d,V}\big(\AA^\bullet,U(B)\big)}{\times} \widetilde{\bord}_{d,V}\big(\AA^{\bullet}\times \AA^1,U(B)\big) \underset{\widetilde{\bord}_{d,V}\big(\AA^\bullet,U(B)\big)}{\times} \widetilde{\bord}_{d,V}\big(\AA^{\bullet},U(B,C)\big).
    \end{equation}
    
    \medskip

    \nid As observed in \Cref{ch2:: obs: discreteness}, weak equivalences between discrete simplicial spaces amount to levelwise set bijections. To show that \eqref{eq:iterated fiber product for Segal condition 1} is weakly equivalent to $\widetilde{\bord}_{d,V}\big(\AA^\bullet,U(A,B,C)\big)$, we adapt the proof of \cite[Theorem~5.4]{BBEP24} to our setting, using the naive sheaf condition from \Cref{ch3:: prop: naive sheaf con}. Concretely, given an open cover $\{U,V\}$ of $[0,1]$, the simplicial set $\bord_d(\AA^\bullet,T)$ is the strict pullback of the diagram
    \begin{equation}
        \bord_d(\AA^\bullet,T)_{|U} \longrightarrow \bord_d(\AA^\bullet,T)_{|U\cap V} \longleftarrow \bord_d(\AA^\bullet,T)_{|V},
    \end{equation} 
    where the vertical subscript bars denote the levelwise restriction of the simplicial set to the corresponding open subset of the interval $[0,1]$. 

    With this in mind, let us pick a cover $\cU = \{U,V\}$ of $[0,1]$ such that $U$ covers $(A\cup B)$ and $V$ covers $(B\cup C)$, which we depict as
    \begin{equation}
        \tikzfig{figures/Ch3/BOR9}{\begin{tikzpicture}
	\begin{pgfonlayer}{nodelayer}
		\node [style=endpoint] (23) at (6, 0) {};
		\node [style=none] (0) at (-6, 0) {};
		\node [style=none] (1) at (6.5, 0) {};
		\node [style=none] (2) at (-4, 0) {};
		\node [style=none] (3) at (4.5, 0) {};
		\node [style=none] (5) at (-2, 0) {};
		\node [style=none] (6) at (-0.5, 0) {};
		\node [style=none] (7) at (0.75, 0) {};
		\node [style=none] (8) at (2.5, 0) {};
		\node [style=none] (9) at (-3, -1) {$A$};
		\node [style=none] (10) at (0.125, -1) {$B$};
		\node [style=none] (11) at (3.5, -1) {$C$};
		\node [style=endpoint] (12) at (-6, 0) {};
		\node [style=none] (14) at (-6, -0.25) {};
		\node [style=none] (15) at (1.25, -0.25) {};
		\node [style=none] (16) at (-6, 0.25) {};
		\node [style=none] (17) at (1.25, 0.25) {};
		\node [style=none] (18) at (-1.25, 0.45) {};
		\node [style=none] (19) at (-1.25, -0.45) {};
		\node [style=none] (20) at (6, 0.45) {};
		\node [style=none] (21) at (6, -0.45) {};
		\node [style=none] (22) at (6, 0) {};
	\end{pgfonlayer}
	\begin{pgfonlayer}{edgelayer}
		\draw [style=dir] (0.center) to (1.center);
		\draw [style={open_int}] (3.center) to (22.center);
		\draw [style={open_int}] (0.center) to (2.center);
		\draw [style={open_int}] (5.center) to (6.center);
		\draw [style={open_int}] (7.center) to (8.center);
		\draw [style=red shade] (17.center)
			 to (16.center)
			 to (14.center)
			 to (15.center)
			 to cycle;
		\draw [style=blue shade] (19.center)
			 to (21.center)
			 to (20.center)
			 to (18.center)
			 to cycle;
		\draw [style=red shade] (3.center) to (22.center);
		\draw [style=red shade] (3.center) to (22.center);
		\draw [style=red shade] (8.center) to (3.center);
	\end{pgfonlayer}
\end{tikzpicture}
}
    \end{equation}
    such that the red shaded region and the blue shaded region cover the opens $U$ and $V$, respectively. Additionally, we define a partition of unity $\{f_{U},f_{V}\}$ subordinate to $\cU$. Define smaller opens $U' = f_{U}^{-1}[0,1/3)$ and $V' = f_V^{-1}(2/3,1]$, noting that, since $f_U + f_V \equiv 1$, they do not overlap, $U' \cap V' = \emptyset$, and thus the cover $\cU' = \{U', U\cap V, V'\}$ refines $\cU$. There is an induced map of simplicial sets
    \begin{equation}
    \begin{tikzcd}
        \bord_d\big(\AA^{\bullet},U(A,B)\big)     
        \underset{\bord_d (\AA^\bullet,U(B))}{\times^{h}} 
        \bord_d\big(\AA^{\bullet}\times \AA^1,U(B)\big) 
        \underset{\bord_d (\AA^\bullet,U(B))}{\times^{h}} 
        \bord_d\big(\AA^{\bullet},U(B,C)\big) \ar[dd, "\mathsf{res}"] \\ \\
        \bord_d\big(\AA^{\bullet},U(A,B)\big)_{|U'}     
        \underset{\bord_d (\AA^\bullet,U(B))_{|U'}}{\times^{h}} 
        \bord_d\big(\AA^{\bullet}\times \AA^1,U(B)\big) 
        \underset{\bord_d (\AA^\bullet,U(B))_{|V'}}{\times^{h}} 
        \bord_d\big(\AA^{\bullet},U(B,C)\big)_{|V'} 
    \end{tikzcd}
    \end{equation}
    which is a weak equivalence, by applying \Cref{ch3:: prop: naive sheaf con} to the covers $\{U',U\cap V\}$ of $U$ and $\{V',U\cap V\}$ of $V$.

    We now construct the following additional pair of simplicial maps:
    \begin{equation}
    \begin{tikzcd}
        \bord_d\big(\AA^{\bullet},U(A,B)\big)     
        \underset{\bord_d\big(\AA^\bullet,U(B)\big)}{\times} 
        \bord_d\big(\AA^{\bullet}\times \AA^1,U(B)\big) 
        \underset{\bord_d\big(\AA^\bullet,U(B)\big)}{\times} 
        \bord_d\big(\AA^{\bullet},U(B,C)\big) \ar[dd, "\mathsf{res}"] \\ \\
        \bord_d\big(\AA^{\bullet},U(A,B)\big)_{|U'}     
        \underset{\bord_d\big(\AA^\bullet,U(B)\big)_{|U'}}{\times} 
        \bord_d\big(\AA^{\bullet}\times \AA^1,U(B)\big) 
        \underset{\bord_d\big(\AA^\bullet,U(B)\big)_{|V'}}{\times} 
        \bord_d\big(\AA^{\bullet},U(B,C)\big)_{|V'} \ar[dd, bend right=25, "F"', shorten <=2pt, shorten >=2pt]   \\ \\ 
        \bord_d\big(\AA^{\bullet},U(A,B)\big)_{|U'}     
        \underset{\bord_d\big(\AA^\bullet,U(B)\big)_{|U'}}{\times} 
        \bord_d\big(\AA^{\bullet},U(B)\big) 
        \underset{\bord_d\big(\AA^\bullet,U(B)\big)_{|V'}}{\times} 
        \bord_d\big(\AA^{\bullet},U(B,C)\big)_{|V'} \ar[uu, bend right=25, "\pr^\ast"']      
    \end{tikzcd}
    \end{equation}
    about which we now expand:
    \begin{enumerate}
        \item As mentioned above, the map $\mathsf{res}$ is an isomorphism of simplicial sets, by \Cref{ch3:: prop: naive sheaf con}.
        \item Crucially, the induced map from $\bord_d\big(\AA^\bullet, U(A,B,C)\big)$ to the fiber product in the bottom of the diagram is also an isomorphism of simplicial sets, again using \Cref{ch3:: prop: naive sheaf con} with the refined cover $\{U',U\cap V,V'\}$.
        \item Lastly, we must articulate why $F$ is a homotopy inverse to $\pr^\ast$. There is a smooth map $c\maps \AA^1 \longra \AA^1$ which is strictly increasing and satisfies the following additional conditions: 
        \begin{enumerate}
        \item $c_{|(-\oo,1/3)} \equiv 0$, and 
        \item $c_{|(2/3,+\oo)} \equiv 1$. 
        \end{enumerate}
        Composing $c \circ (f_{V})_{|U\cap V}\maps U\cap V \longra \AA^1$ we get a function that vanishes on $U'\subset U\cap V$ and is constantly equal to $1$ on $V'\subset U\cap V$; we call this map $f$, which we depict in the following blue graph, as an aid to the reader:
        \begin{equation}
            \scalebox{0.8}{\tikzfig{figures/Ch3/BOR10}{\begin{tikzpicture}
	\begin{pgfonlayer}{nodelayer}
		\node [style=endpoint] (23) at (6, -0.75) {};
		\node [style=none] (0) at (-6, -0.75) {};
		\node [style=none] (1) at (6.5, -0.75) {};
		\node [style=none] (2) at (-4, -0.75) {};
		\node [style=none] (3) at (4.5, -0.75) {};
		\node [style=none] (5) at (-2.25, -0.75) {};
		\node [style=none] (6) at (-0.5, -0.75) {};
		\node [style=none] (7) at (0.75, -0.75) {};
		\node [style=none] (8) at (2.5, -0.75) {};
		\node [style=none] (9) at (-3, -1.75) {$A$};
		\node [style=none] (10) at (0.125, -1.75) {$B$};
		\node [style=none] (11) at (3.5, -1.75) {$C$};
		\node [style=endpoint] (12) at (-6, -0.75) {};
		\node [style=none] (14) at (-6, -1) {};
		\node [style=none] (15) at (1.5, -1) {};
		\node [style=none] (16) at (-6, -0.5) {};
		\node [style=none] (17) at (1.5, -0.5) {};
		\node [style=none] (18) at (-2, -0.3) {};
		\node [style=none] (19) at (-2, -1.2) {};
		\node [style=none] (20) at (6, -0.3) {};
		\node [style=none] (21) at (6, -1.2) {};
		\node [style=none] (22) at (6, -0.75) {};
		\node [style=none] (27) at (-1.25, -0.75) {};
		\node [style=none] (28) at (1.25, -0.75) {};
		\node [style=none] (29) at (-6, 0) {};
		\node [style=none] (30) at (6.5, 0) {};
		\node [style=none] (31) at (-6, 9.75) {};
		\node [style=none] (32) at (-1.75, 9) {};
		\node [style=none] (33) at (-6, 9) {};
		\node [style=none] (34) at (1.5, 0.025) {};
		\node [style=none] (35) at (6, 0) {};
		\node [style=none] (36) at (6, 9) {};
		\node [style=none] (37) at (-2, 0) {};
		\node [style=none] (38) at (1.5, 9) {};
		\node [style=none] (39) at (4, 9.5) {$f_{V}$};
		\node [style=none] (40) at (-3.75, 9.5) {$f_{U}$};
		\node [style=none] (41) at (-6, 6) {};
		\node [style=none] (42) at (-6, 3) {};
		\node [style=none] (43) at (-7, 3) {$1/3$};
		\node [style=none] (44) at (-7, 6) {$2/3$};
		\node [style=none] (45) at (-6.25, 3) {};
		\node [style=none] (46) at (-5.75, 3) {};
		\node [style=none] (47) at (-6.25, 6) {};
		\node [style=none] (48) at (-5.75, 6) {};
		\node [style=none] (49) at (6, 6) {};
		\node [style={endpoint_blue}] (50) at (0.05, 6.025) {};
		\node [style=none] (51) at (-0.625, 0) {};
		\node [style=none] (52) at (6, 3) {};
		\node [style={endpoint_blue}] (53) at (-0.525, 3) {};
		\node [style=none] (54) at (0.25, 9) {};
		\node [style=none] (55) at (-6, 10.5) {$\AA^1$};
		\node [style=none] (56) at (-7, 9) {$1$};
		
	\end{pgfonlayer}
	\begin{pgfonlayer}{edgelayer}
		\draw [style=dir] (0.center) to (1.center);
		\draw [style={open_int}] (3.center) to (22.center);
		\draw [style={open_int}] (0.center) to (2.center);
		\draw [style={open_int}] (5.center) to (6.center);
		\draw [style={open_int}] (7.center) to (8.center);
		\draw [style=red shade] (17.center)
			 to (16.center)
			 to (14.center)
			 to (15.center)
			 to cycle;
		\draw [style=blue shade] (19.center)
			 to (21.center)
			 to (20.center)
			 to (18.center)
			 to cycle;
		\draw [style=red shade] (3.center) to (22.center);
		\draw [style=red shade] (3.center) to (22.center);
		\draw [style=red shade] (8.center) to (3.center);
		\draw (7.center) to (3.center);
		\draw [style=dir] (29.center) to (30.center);
		\draw [style=dir] (29.center) to (31.center);
		\draw [style=geometric] (33.center) to (32.center);
		\draw [style=geometric, in=180, out=0, looseness=0.50] (32.center) to (34.center);
		\draw [style=geometric] (34.center) to (35.center);
		\draw [style=geometric] (36.center) to (38.center);
		\draw [style=geometric, in=0, out=-180, looseness=0.50] (38.center) to (37.center);
		\draw [style=geometric] (37.center) to (29.center);
		\draw (45.center) to (46.center);
		\draw (47.center) to (48.center);
		\draw [style=ts-dashed] (48.center) to (49.center);
		\draw [style={open_int_blue}] (37.center) to (51.center);
		\draw [style=ts-dashed] (46.center) to (52.center);
		\draw [style={open_int_blue}, in=-180, out=0, looseness=0.25] (51.center) to (54.center);
		\draw [style={open_int_blue}] (54.center) to (38.center);
	\end{pgfonlayer}
\end{tikzpicture}
}}
        \end{equation}
        The map $F$ sends a smooth submanifold 
        \[
        W\subset (\AA^\bullet \times \AA^1)\times U(B)\times V
        \]
        to the pullback along 
        \begin{equation}
        \id_{\AA^{\bullet}}\times f \times \id_{V}\maps \AA^{\bullet}\times U(B) \times V \longra \AA^{\bullet}\times \AA^1 \times U(B) \times V,
        \end{equation}
        obtaining a smooth $(d+\bullet + 1)$-dimensional submanifold $\big(W_f \subset \AA^\bullet\times U(B)\times V\big)$. For intuition, we sketch its action in dimension $d=1$:
        \begin{equation}
            \!\!\!\!\scalebox{0.8}{\tikzfig{figures/Ch3/BOR11}{\begin{tikzpicture}
	\begin{pgfonlayer}{nodelayer}
		\node [style=endpoint] (23) at (9, 3) {};
		\node [style=none] (0) at (-3, 3) {};
		\node [style=none] (1) at (9.5, 3) {};
		\node [style=none] (2) at (-1, 3) {};
		\node [style=none] (3) at (7.5, 3) {};
		\node [style=none] (5) at (-0.5, 3) {};
		\node [style=none] (6) at (0.75, 3) {};
		\node [style=none] (7) at (5.25, 3) {};
		\node [style=none] (8) at (6.75, 3) {};
		\node [style=endpoint] (12) at (-3, 3) {};
		\node [style=none] (18) at (1.775, 3) {};
		\node [style=none] (20) at (9, 3) {};
		\node [style=none] (21) at (4, -5) {};
		\node [style=none] (22) at (9, 3) {};
		\node [style=endpoint] (24) at (4, -5) {};
		\node [style=none] (25) at (-7.95, -5) {};
		\node [style=none] (26) at (4.5, -5) {};
		\node [style=none] (27) at (-6, -5) {};
		\node [style=none] (28) at (2.5, -5) {};
		\node [style=none] (29) at (-5.5, -5) {};
		\node [style=none] (30) at (-4.25, -5) {};
		\node [style=none] (31) at (0.25, -5) {};
		\node [style=none] (32) at (1.75, -5) {};
		\node [style=none] (33) at (-5.75, -5.75) {$A$};
		\node [style=none] (34) at (-2.125, -5.75) {$B$};
		\node [style=none] (35) at (2, -5.75) {$C$};
		\node [style=none] (45) at (4, -5) {};
		\node [style=none] (46) at (-3, 10) {};
		\node [style=none] (47) at (-3, 3) {};
		\node [style=none] (48) at (-8.25, -5.5) {};
		\node [style=none] (49) at (-3, 3) {};
		\node [style=none] (71) at (3.75, -5.5) {};
		\node [style=none] (72) at (9, 3) {};
		\node [style=none] (73) at (4.5, 3) {};
		\node [style=none] (75) at (-0.25, 3) {};
		\node [style=none] (76) at (6.25, 3) {};
		\node [style=none] (77) at (1.25, -5) {};
		\node [style=none] (78) at (-3, 3) {};
		\node [style=none] (79) at (6.25, 3) {};
		\node [style=none] (80) at (-8, -5) {};
		\node [style=none] (81) at (1.25, -5) {};
		\node [style=none] (82) at (-0.25, 3) {};
		\node [style=none] (83) at (9, 3) {};
		\node [style=none] (84) at (-5.25, -5) {};
		\node [style=none] (85) at (4, -5) {};
		\node [style=none] (50) at (0.75, 6) {};
		\node [style=none] (51) at (5.25, 8) {};
		\node [style=none] (69) at (-4.25, -2) {};
		\node [style=none] (70) at (0.25, 0.75) {};
		\node [style=none] (86) at (-4.5, 10) {$\AA^\bullet \times V$};
		\node [style=none] (87) at (-9, -5.5) {$\AA^1$};
	\end{pgfonlayer}
	\begin{pgfonlayer}{edgelayer}
		\draw [style=ts-dashed] (50.center) to (6.center);
		\draw [style=dir] (0.center) to (1.center);
		\draw [style={open_int}] (3.center) to (22.center);
		\draw [style={open_int}] (0.center) to (2.center);
		\draw [style={open_int}] (5.center) to (6.center);
		\draw [style={open_int}] (7.center) to (8.center);
		\draw [style=blue shade] (21.center) to (20.center);
		\draw [style=blue shade] (20.center) to (18.center);
		\draw [style=red shade] (3.center) to (22.center);
		\draw [style=red shade] (3.center) to (22.center);
		\draw [style=red shade] (8.center) to (3.center);
		\draw [style=dir] (25.center) to (26.center);
		\draw [style={open_int}] (28.center) to (45.center);
		\draw [style={open_int}] (25.center) to (27.center);
		\draw [style={open_int}] (29.center) to (30.center);
		\draw [style={open_int}] (31.center) to (32.center);
		\draw [style=red shade] (28.center) to (45.center);
		\draw [style=red shade] (28.center) to (45.center);
		\draw [style=red shade] (32.center) to (28.center);
		\draw [style=dir] (47.center) to (48.center);
		\draw [style=dir] (49.center) to (46.center);
		\draw [style=red shade] (77.center) to (76.center);
		\draw [style=red shade] (76.center) to (49.center);
		\draw [style=red shade, in=135, out=45, loop] (49.center) to ();
		\draw [style=red shade] (49.center) to (73.center);
		\draw [style=blue shade] (45.center) to (72.center);
		\draw [style=blue shade] (72.center) to (75.center);
		\draw [style=red shade] (76.center) to (77.center);
		\draw [style=red shade] (77.center) to (25.center);
		\draw [style=red shade] (25.center) to (49.center);
		\draw [style=red shade] (49.center) to (76.center);
		\draw [style=red shade] (80.center)
			 to (78.center)
			 to (79.center)
			 to (81.center)
			 to cycle;
		\draw [style=blue shade] (85.center)
			 to (84.center)
			 to (82.center)
			 to (83.center)
			 to cycle;
		\draw [style=geometric shaded] (69.center)
			 to [in=165, out=0] (70.center)
			 to (51.center)
			 to [in=0, out=-180] (50.center)
			 to [in=60, out=-120] cycle;
		\draw [style=ts-dashed] (70.center) to (31.center);
		\draw [style=ts-dashed] (69.center) to (30.center);
		\draw [style=ts-dashed] (51.center) to (7.center);
	\end{pgfonlayer}
\end{tikzpicture}
}} \rightsquigarrow \scalebox{0.75}{\tikzfig{figures/Ch3/BOR12}{\begin{tikzpicture}
	\begin{pgfonlayer}{nodelayer}
		\node [style=endpoint] (23) at (9, 3) {};
		\node [style=none] (0) at (-3, 3) {};
		\node [style=none] (1) at (9.5, 3) {};
		\node [style=none] (2) at (-1, 3) {};
		\node [style=none] (3) at (8, 3) {};
		\node [style=none] (5) at (-0.5, 3) {};
		\node [style=none] (6) at (0.75, 3) {};
		\node [style=none] (7) at (5.25, 3) {};
		\node [style=none] (8) at (6.75, 3) {};
		\node [style=endpoint] (12) at (-3, 3) {};
		\node [style=none] (18) at (1.775, 3) {};
		\node [style=none] (20) at (9, 3) {};
		\node [style=none] (21) at (4, -5) {};
		\node [style=none] (22) at (9, 3) {};
		\node [style=endpoint] (24) at (4, -5) {};
		\node [style=none] (25) at (-7.95, -5) {};
		\node [style=none] (26) at (4.5, -5) {};
		\node [style=none] (27) at (-6, -5) {};
		\node [style=none] (28) at (3, -5) {};
		\node [style=none] (29) at (-5.5, -5) {};
		\node [style=none] (30) at (-4.25, -5) {};
		\node [style=none] (31) at (1.25, -5) {};
		\node [style=none] (32) at (1.75, -5) {};
		\node [style=none] (33) at (-5.75, -5.75) {$A$};
		\node [style=none] (34) at (-2.125, -5.75) {$B$};
		\node [style=none] (35) at (2, -5.75) {$C$};
		\node [style=none] (45) at (4, -5) {};
		\node [style=none] (46) at (-3, 10) {};
		\node [style=none] (47) at (-3, 3) {};
		\node [style=none] (48) at (-8.25, -5.5) {};
		\node [style=none] (49) at (-3, 3) {};
		\node [style=none] (71) at (3.75, -5.5) {};
		\node [style=none] (72) at (9, 3) {};
		\node [style=none] (73) at (4.5, 3) {};
		\node [style=none] (75) at (-0.25, 3) {};
		\node [style=none] (76) at (6.25, 3) {};
		\node [style=none] (77) at (1.25, -5) {};
		\node [style=none] (86) at (-4.5, 10) {$\AA^\bullet \times V$};
		\node [style=none] (87) at (-9, -5.5) {$\AA^1$};
		\node [style=none] (88) at (0.75, 3) {};
		\node [style=none] (89) at (2, 3) {};
		\node [style=none] (90) at (-0.5, -5) {};
		\node [style=none] (91) at (1.25, -5) {};
		\node [style=none] (92) at (2, 6.3) {};
		\node [style=none] (50) at (0.75, 6) {};
		\node [style=none] (51) at (5.25, 8) {};
		\node [style=none] (69) at (-4.25, -2) {};
		\node [style=none] (70) at (1.25, -4.5) {};
	\end{pgfonlayer}
	\begin{pgfonlayer}{edgelayer}
		\draw [style=ts-dashed] (50.center) to (6.center);
		\draw [style=dir] (0.center) to (1.center);
		\draw [style={open_int}] (3.center) to (22.center);
		\draw [style={open_int}] (0.center) to (2.center);
		\draw [style={open_int}] (5.center) to (6.center);
		\draw [style={open_int}] (7.center) to (8.center);
		\draw [style=blue shade] (21.center) to (20.center);
		\draw [style=blue shade] (20.center) to (18.center);
		\draw [style=red shade] (3.center) to (22.center);
		\draw [style=red shade] (3.center) to (22.center);
		\draw [style=red shade] (8.center) to (3.center);
		\draw [style=dir] (25.center) to (26.center);
		\draw [style={open_int}] (28.center) to (45.center);
		\draw [style={open_int}] (25.center) to (27.center);
		\draw [style={open_int}] (29.center) to (30.center);
		\draw [style={open_int}] (31.center) to (32.center);
		\draw [style=red shade] (28.center) to (45.center);
		\draw [style=red shade] (28.center) to (45.center);
		\draw [style=red shade] (32.center) to (28.center);
		\draw [style=dir] (47.center) to (48.center);
		\draw [style=dir] (49.center) to (46.center);
		\draw [style=red shade] (77.center) to (76.center);
		\draw [style=red shade] (76.center) to (49.center);
		\draw [style=red shade, in=135, out=45, loop] (49.center) to ();
		\draw [style=red shade] (49.center) to (73.center);
		\draw [style=blue shade] (45.center) to (72.center);
		\draw [style=blue shade] (72.center) to (75.center);
		\draw [style=red shade] (76.center) to (77.center);
		\draw [style=red shade] (77.center) to (25.center);
		\draw [style=red shade] (25.center) to (49.center);
		\draw [style=red shade] (49.center) to (76.center);
		\draw [style=ts-dashed] (69.center) to (30.center);
		\draw [style=ts-dashed] (51.center) to (7.center);
		\draw [style=ts-dashed, in=90, out=-90] (70.center) to (31.center);
		\draw [style=geometric shaded] (69.center)
			 to [in=165, out=0] (70.center)
			 to (51.center)
			 to [in=0, out=-180] (50.center)
			 to [in=60, out=-120] cycle;
	\end{pgfonlayer}
\end{tikzpicture}
}}
        \end{equation}
        On the left we drew a section, i.e.\ a $(d+1+\bullet)$-dimensional manifold admitting a submersion onto $\AA^{1}\times U(B)$. On the right we depict the interpolating submanifold induced by $f$. Lastly, the next figure captures $(\id_{\AA^\bullet} \times f \times \id_V)$ with the orange graph; we leave the fibers over $U' \cap U(B)$ and $U(B)\cap V'$ in the original green so as to highlight how $W_f$ interpolates the two:
        \begin{equation}
            \scalebox{0.8}{\tikzfig{figures/Ch3/BOR13}{\begin{tikzpicture}
	\begin{pgfonlayer}{nodelayer}
		\node [style=endpoint] (23) at (9, 3) {};
		\node [style=none] (0) at (-3, 3) {};
		\node [style=none] (1) at (9.5, 3) {};
		\node [style=none] (2) at (-1, 3) {};
		\node [style=none] (3) at (7.5, 3) {};
		\node [style=none] (5) at (-0.5, 3) {};
		\node [style=none] (6) at (0.75, 3) {};
		\node [style=none] (7) at (5.25, 3) {};
		\node [style=none] (8) at (6.75, 3) {};
		\node [style=endpoint] (12) at (-3, 3) {};
		\node [style=none] (18) at (1.775, 3) {};
		\node [style=none] (20) at (9, 3) {};
		\node [style=none] (21) at (4, -5) {};
		\node [style=none] (22) at (9, 3) {};
		\node [style=endpoint] (24) at (4, -5) {};
		\node [style=none] (25) at (-7.95, -5) {};
		\node [style=none] (26) at (4.5, -5) {};
		\node [style=none] (27) at (-6, -5) {};
		\node [style=none] (28) at (2.5, -5) {};
		\node [style=none] (29) at (-5.5, -5) {};
		\node [style=none] (30) at (-4.25, -5) {};
		\node [style=none] (31) at (1.25, -5) {};
		\node [style=none] (32) at (1.75, -5) {};
		\node [style=none] (33) at (-5.75, -5.75) {$A$};
		\node [style=none] (34) at (-2.125, -5.75) {$B$};
		\node [style=none] (35) at (2, -5.75) {$C$};
		\node [style=none] (45) at (4, -5) {};
		\node [style=none] (46) at (-3, 10) {};
		\node [style=none] (47) at (-3, 3) {};
		\node [style=none] (48) at (-8.25, -5.5) {};
		\node [style=none] (49) at (-3, 3) {};
		\node [style=none] (71) at (3.75, -5.5) {};
		\node [style=none] (72) at (9, 3) {};
		\node [style=none] (73) at (4.5, 3) {};
		\node [style=none] (75) at (-0.25, 3) {};
		\node [style=none] (76) at (6.25, 3) {};
		\node [style=none] (77) at (1.25, -5) {};
		\node [style=none] (78) at (-3, 3) {};
		\node [style=none] (79) at (6.25, 3) {};
		\node [style=none] (80) at (-8, -5) {};
		\node [style=none] (81) at (1.25, -5) {};
		\node [style=none] (82) at (-0.25, 3) {};
		\node [style=none] (83) at (9, 3) {};
		\node [style=none] (84) at (-5.25, -5) {};
		\node [style=none] (85) at (4, -5) {};
		\node [style=none] (86) at (-4, 10) {$\scriptstyle\AA^\bullet \times V$};
		\node [style=none] (87) at (-9, -5.5) {$\scriptstyle\AA^1$};
		\node [style=none] (88) at (0.75, 3) {};
		\node [style=none] (89) at (2, 3) {};
		\node [style=none] (90) at (-0.5, -5) {};
		\node [style=none] (91) at (1.25, -5) {};
		\node [style=none] (92) at (2, 6.25) {};
		\node [style=none] (93) at (-0.5, 0.45) {};
		\node [style=none] (50) at (0.75, 6) {};
		\node [style=none] (51) at (5.25, 7.25) {};
		\node [style=none] (69) at (-4.25, -2) {};
		\node [style=none] (70) at (1.25, 0.75) {};
		\node [style=none] (94) at (3.5, 8.2) {};
		\node [style=none] (95) at (5.25, 8.5) {};
		\node [style=none] (96) at (2.5, 8.25) {$W_f$};
	\end{pgfonlayer}
	\begin{pgfonlayer}{edgelayer}
		\draw [style=ts-dashed] (50.center) to (6.center);
		\draw [style=dir] (0.center) to (1.center);
		\draw [style={open_int}] (3.center) to (22.center);
		\draw [style={open_int}] (0.center) to (2.center);
		\draw [style={open_int}] (5.center) to (6.center);
		\draw [style={open_int}] (7.center) to (8.center);
		\draw [style=blue shade] (21.center) to (20.center);
		\draw [style=blue shade] (20.center) to (18.center);
		\draw [style=red shade] (3.center) to (22.center);
		\draw [style=red shade] (3.center) to (22.center);
		\draw [style=red shade] (8.center) to (3.center);
		\draw [style=dir] (25.center) to (26.center);
		\draw [style={open_int}] (28.center) to (45.center);
		\draw [style={open_int}] (25.center) to (27.center);
		\draw [style={open_int}] (29.center) to (30.center);
		\draw [style={open_int}] (31.center) to (32.center);
		\draw [style=red shade] (28.center) to (45.center);
		\draw [style=red shade] (28.center) to (45.center);
		\draw [style=red shade] (32.center) to (28.center);
		\draw [style=dir] (47.center) to (48.center);
		\draw [style=dir] (49.center) to (46.center);
		\draw [style=red shade] (77.center) to (76.center);
		\draw [style=red shade] (76.center) to (49.center);
		\draw [style=red shade, in=135, out=45, loop] (49.center) to ();
		\draw [style=red shade] (49.center) to (73.center);
		\draw [style=blue shade] (45.center) to (72.center);
		\draw [style=blue shade] (72.center) to (75.center);
		\draw [style=red shade] (76.center) to (77.center);
		\draw [style=red shade] (77.center) to (25.center);
		\draw [style=red shade] (25.center) to (49.center);
		\draw [style=red shade] (49.center) to (76.center);
		\draw [style=red shade] (80.center)
			 to (78.center)
			 to (79.center)
			 to (81.center)
			 to cycle;
		\draw [style=blue shade] (85.center)
			 to (84.center)
			 to (82.center)
			 to (83.center)
			 to cycle;
		\draw [style=ts-dashed] (69.center) to (30.center);
		\draw [style=ts-dashed] (51.center) to (7.center);
		\draw [style={open_int_blue}] (88.center) to (89.center);
		\draw [style={open_int_blue}] (90.center) to (91.center);
		\draw [style={open_int_blue}, in=-180, out=0, looseness=0.25] (89.center) to (90.center);
		\draw [style=ts-dashed, in=90, out=-90] (70.center) to (31.center);
		\draw [style=geometric shaded] (69.center)
			 to [in=165, out=0] (70.center)
			 to (51.center)
			 to [in=0, out=-180] (50.center)
			 to [in=60, out=-120] cycle;
		\draw [style={open_int_green}, in=-135, out=30, looseness=0.50] (92.center) to (93.center);
		\draw [style={open_int_green}, in=165, out=45, looseness=0.75] (93.center) to (70.center);
		\draw [style={open_int_green}, in=-150, out=0, looseness=0.50] (50.center) to (92.center);
		\draw [style={open_int_green}, in=165, out=45, looseness=0.75] (94.center) to (95.center);
		\draw [style=new edge style 0, in=225, out=30, looseness=1.25] (92.center) to (94.center);
	\end{pgfonlayer}
\end{tikzpicture}
}}
        \end{equation}
        This figure is meant to elucidate this map's well-definedness --- concretely, the fact that it preserves the source at $U'$ and the target at $V'$:
         \begin{equation}
        (W_f)_{|U'} = (\mathsf{source})^\ast(W)_{|U'}, \quad
        (W_f)_{|V'} = (\mathsf{target})^\ast(W)_{|V'}, 
        \end{equation} 
        which is a direct consequence of how $f$ is defined as a constant map on $0$ and $1$ over $U'$ and $V'$, respectively. Observe that $(\pr^\ast W)_f = (W \times \AA^1)_f =  W$, stemming from the fact that 
        \begin{equation}
            \AA^\bullet \times U(B)\times V \xra{\,\id \times f \times \id\,} \AA^\bullet \times \AA^1 \times U(B)\times V \xra{\,\pr\,} \AA^\bullet \times U(B) \times V = \id,
        \end{equation}
        thus proving that $F \circ \pr^\ast = \id$. Intuitively, this map sends a $d$-dimensional submanifold $W$ to its corresponding cylinder over $\AA^1$ and then applies $F$, recovering $W$. 
        \item Finally, we must prove that $\pr^* \circ F \simeq \id$. The composite $\pr^* \circ F$ sends a submanifold 
        \begin{equation}
            \big(W \subset (\AA^\bullet \times \AA^1) \times U(B) \times V \big) \longmapsto \big((W_f \times \AA^1) \subset (\AA^\bullet \times \AA^1) \times U(B) \times V \big).
        \end{equation} 
        We are interested in finding a simplicial homotopy 
        \begin{equation}
            \eta \maps \bord_d(\AA^\bullet \times \AA^{1}, U(B)) \times \Delta^1 \longra \bord_d(\AA^\bullet \times \AA^1 , U(B))
        \end{equation}
        from $\pr^\ast \circ F$ to $\id$. 
        We first apply the adjunction to reframe the homotopy as a simplicial map to the path space  
        \begin{equation}
        \begin{tikzcd}[ampersand replacement=\&]
            {\bord_d(\AA^{\bullet}\times \AA^{1},U(B))} \&\& {\mathsf{Map}_{\sset}\big(\Delta^1,\bord_d(\AA^{\bullet}\times \AA^{1},U(B))\big)} \\
            \\
            \&\& {\bord_d\big((\AA^{\bullet}\times \AA^{1})\times \AA^1,U(B)\big)}
            \arrow["\eta", from=1-1, to=1-3]
            \arrow["(*)"', curve={height=18pt}, from=1-1, to=3-3]
            \arrow["\simeq", from=1-3, to=3-3]
        \end{tikzcd}
        \end{equation}
        and then, using concordance-invariance, we get that the homotopy is ultimately equivalent to a simplicial map to $\bord_d\big((\AA^{\bullet}\times \AA^{1})\times \AA^1,U(B)\big)$. Let us consider the smooth map
        \begin{equation}
            \id_{\AA^\bullet} \times \,H \times \id_V \maps \big((\AA^\bullet \times \AA^1) \times \AA^1 \times U(B) \times V \big) \longra \big((\AA^\bullet \times \AA^1)  \times U(B) \times V \big) 
        \end{equation}
        where $H\maps (U(B)\times \AA^1) \times \AA^1 \longra U(B) \times \AA^1$ is defined as the smooth homotopy 
        \begin{equation}
            H\big((x,t),s\big) = \big(x, h\big((x,t),s)\big)
        \end{equation}
        via the smooth linear interpolation
        \begin{align}
        h\big((x,t),s\big) &= (1-s)\cdot f \circ \pr_{U(B)}(x,t) + s \cdot \pr_{\AA^1}(x,t) \\
        &= (1-s)\cdot f(x) + s\cdot t.
        \end{align}
        The smooth homotopy $H$ starts at $H\big((x,t),0\big) = (x,f(x))$ and ends at $H\big((x,t),1\big) = (x,t)$, i.e.\ the identity. Now we check that the diagonal map $(*)$ satisfies the properties we are after: when we ``evaluate'' the homotopy at $\{i\}$, by pulling back along $\AA^0 \xra{\langle i \rangle^*} \AA^1$, we have that:
        \begin{equation}
            \big(\id\times H \times \id\big)((u,i),t,x,v) = \begin{cases}
                (u,f(x),x,v) & \text{ if } i=0, \\
                (u,t,x,v) & \text{ if } i = 1.
            \end{cases}
        \end{equation}
        For $i=1$ we manifestly get the identity, as desired, and for $i=0$ we have 
        \begin{equation}
        (\id \times \,\pr\times \id) \circ(\id_{\AA^\bullet}\times \,f\times \id_V),
        \end{equation}
        thus finishing the proof.  
    \end{enumerate}
\end{proof}




\begin{observation}
We pause here to reflect on the difference between this proof of the Segal condition and that of Calaque--Scheimbauer \cite[Proposition 5.19]{CS19}. Note first that we dispense with the need to pass through various subspaces of the homotopy fiber product, and, more importantly, we get to use the intrinsic geometric nature of our construction as the conduit for simplicial homotopy theory. This feature, we believe, may pay dividends in further applications to manifold topology and topological field theories. In a way, we see our construction as a means to discretize their $(\oo,1)$-category of intervals, thus simplifying it further.
\end{observation}


The following result characterizes the structural conditions that any model of the bordism category must satisfy. To make sense of these, we must first introduce some less well-studied Lie groups first. As far as we are aware, they were first defined in \cite{GMTW09}. 

\begin{definition}\label{ch3:: def: iso bordism}
 Let $(M,h_0,h_1)$ be a $d$-dimensional smooth unoriented bordism from $\Sigma_0$ to $\Sigma_1$ where we have collars
 \begin{align}
    h_0 \maps \Sigma_0 \times [0,1) &\hookra M \\
    h_1 \maps \Sigma_1 \times (0,1] &\hookra M
 \end{align}
 where $\partial M$ is the image of $h_{i}(\{i\}\times \Sigma_i)$ for $i=0,1$. Let $\Diff(\{M,h_0,h_1\})$ be the subgroup of diffeomorphisms of $M$ which must satisfy the following additional condition: 
 \begin{itemize}
    \item[($*$)] There exists a diffeomorphism $[0,1) \longra [0,1)$ making the following diagram commute:
    \begin{equation}
        \begin{tikzcd}
            {h_0(\Sigma_0\times [0,1))} && {\varphi\big(h_0(\Sigma_0\times [0,1))\big)} \\
            \\
            {[0,1)} && {[0,1)}
            \arrow["\varphi", from=1-1, to=1-3]
            \arrow[from=1-1, to=3-1]
            \arrow[from=1-3, to=3-3]
            \arrow[dashed, from=3-1, to=3-3]
        \end{tikzcd}
    \end{equation} 
    with an analogous condition for $h_1$.
 \end{itemize}
\end{definition}

\begin{observation}
    It is instructive to think of the group $\Diff(\{M,h_0,h_1\})$ as the subgroup of diffeomorphisms of $M$ which are a product of diffeomorphisms near the boundary $\partial M$.
\end{observation}

We want to not only consider bordisms, as they are interpreted as 1-morphisms, but also chains of composable bordisms. Moreover, we want to study the space of such composable chains, and so we must, in some fashion, generalize the diffeomorphism group in order to accommodate these additional blocks of data.

\begin{definition}\label{ch3:: def: diffeo group}
    Let $(M,\{h_i\}_{i=0}^{k})$ be a composed chain of $k$-composable $d$-dimensional unoriented bordisms 
    \begin{equation}
        \Sigma_0 \xra{M_1} \Sigma_1 \xra{M_2} \cdots \xra{M_{k}} \Sigma_k 
    \end{equation}
    together with the accompanying collaring data 
    \begin{align}
        h_0 \maps [0,1) \times &\Sigma_0 \hookra M \\
        h_i \maps (0,1) \times &\Sigma_i \hookra M \\
        h_k \maps (0,1]  \times &\Sigma_k \hookra M.
    \end{align}
    Let $\Diff(M,\{h_i\}_{i=0}^{k})$ denote the subgroup of diffeomorphisms of $M$ which satisfy the additional set of conditions:
    \begin{itemize}
        \item[$(*)$] There exist diffeomorphisms $(0,1) \longra (0,1)$ making the following diagram commute:
        \begin{equation}
            \begin{tikzcd}
                {h_i(M_{i}\times(0,1))} && {\varphi\big(h_i(M_{i}\times(0,1))\big)} \\
                \\
                {(0,1)} && {(0,1)}
                \arrow["\varphi", from=1-1, to=1-3]
                \arrow[from=1-1, to=3-1]
                \arrow[from=1-3, to=3-3]
                \arrow[dashed, from=3-1, to=3-3]
            \end{tikzcd}
        \end{equation} 
        with analogous conditions for $h_0$ and $h_1$.
    \end{itemize}
\end{definition}




We now record the fundamental properties an $(\oo,1)$-category $X$ must have in order to be deemed an $(\oo,1)$-category of $d$-dimensional unoriented bordisms: 

\begin{definition}
Let $\cF$ be an $(\oo,1)$-category. We say that $\cF$ \emph{is} the $(\oo,1)$-category of $d$-dimensional bordisms if it satisfies the following three conditions:
\begin{enumerate}
    \item \label{ch3:: cond: bordism cat objects} The underlying space of objects $\Obj(\cF)$ must be weakly equivalent to 
    \begin{equation}
        \coprod_{[\Sigma]}\mathbf{B}\!\Diff(\Sigma),
    \end{equation}
    where the coproduct ranges over all diffeomorphism classes of \emph{closed} $(d-1)$-dimensional unoriented manifolds.  
    \item \label{ch3:: cond: bordism cat morphisms} The space $\Hom(\cF)$ of morphisms should be weakly equivalent to the space 
    \begin{equation}
        \coprod_{[M,\partial_{\nu}M_i]}\mathbf{B}\!\Diff(M,\{\partial_{\nu}M\}),
    \end{equation} 
    where the coproduct ranges over all diffeomorphisms classes of unoriented $d$-dimensional cobordisms. Here $\Diff(M,\{\partial_\nu M\})$ is the subgroup of diffeomorphisms of $M$ that restrict to diffeomorphisms of collars on $\partial_0 M$ and $\partial_1 M$. 
    \item \label{ch3:: cond: bordism cat composable} The space of $k$-composable morphisms
    \begin{equation}
        \underbrace{ \Hom(\cF) \underset{\Obj(\cF)}{\times} \Hom(\cF) \underset{\Obj(\cF)}{\times} \cdots \underset{\Obj(\cF)}{\times} \Hom(\cF) \underset{\Obj(\cF)}{\times} \Hom(\cF) }_{k \text{ times}}
    \end{equation} 
    should be weakly equivalent to the space 
    \begin{equation}
         \coprod_{[M_0 \longra \cdots \longra M_k]}\mathbf{B}\!\Diff(M,\{h_i\}_{i=0}^{k}),
    \end{equation}
    where the coproduct is considered over all the diffeomorphism classes of $k$-composable chains of unoriented $d$-dimensional bordisms. This diffeomorphism group contains all the diffeomorphisms of $M$, the composed bordism, which additionally restrict to diffeomorphisms on a tubular neighborhood of $M_i \hookra M$ for all $i=1,\ldots,k$.
\end{enumerate}

\end{definition}

\begin{observation}
We acknowledge the redundancy in the above definition. Condition~(\ref{ch3:: cond: bordism cat morphisms}) is implied by condition~(\ref{ch3:: cond: bordism cat composable}), yet we chose to separate them for clarity.
\end{observation}

\begin{theorem}\label{ch3:: thm: bordism category characterization}
The $(\oo,1)$-category $\bord_{d}$ is \emph{the} $d$-dimensional bordism category.
\end{theorem}

\begin{proof}
    We verify the three conditions from the previous definition. Since $\bord_d$ is an interval Segal space, its values on any two interval objects with the same image under $\gamma$ are weakly equivalent. Thus, for the three conditions, it is enough to work with fixed representatives in $\cJ_0$, $\cJ_1$, and more generally in $\cJ_k$.

    For the object space, fix $\sU(I)\in \cJ_0$. A $k$-simplex of $\bord_d\big(\sU(I)\big)$ is a smooth embedded submanifold
    \[
        W\subset \AA^k\times \core\big(\sU(I)\big)\times \RR^\infty
    \]
    such that the projection
    \[
        \pi_W \maps W\longra \AA^k\times \core\big(\sU(I)\big)
    \]
    is a proper submersion. By \Cref{ap:: prop: Ehresmann fibration no boundary}, $\pi_W$ is therefore a smooth fiber bundle with closed fiber some $(d-1)$-dimensional manifold $\Sigma$. Since the base $\AA^k\times \core\big(\sU(I)\big)$ is contractible, \Cref{ap:: prop: triviality contractibility} identifies this bundle with the trivial bundle
    \[
        \Sigma\times \AA^k\times \core\big(\sU(I)\big).
    \]
    Write
    \[
        B = \AA^k\times \core\big(\sU(I)\big).
    \]
    Choosing a trivialization of $\pi_W$ identifies $W$ with a smoothly varying family of embedded copies of $\Sigma$ inside $\RR^\infty$, parametrized by $B$, hence with a smooth map
    \[
        B\longra \Emb(\Sigma,\RR^\infty).
    \]
    A second trivialization differs from the first by an automorphism of the trivial bundle $\Sigma\times B\to B$, and therefore induces a smooth map $B\to \Diff(\Sigma)$. Thus the resulting parametrized family of embeddings is well-defined only up to the right action of $\Diff(\Sigma)$. In other words, the connected component of $\bord_d\big(\sU(I)\big)$ corresponding to the diffeomorphism type of $\Sigma$ is the usual simplicial model for the moduli of embeddings of $\Sigma$ into $\RR^\infty$ modulo reparametrization by $\Diff(\Sigma)$. This is precisely the embedded-bundle picture used in \cite[\S2.1, p.~4]{GMTW09}: there the orbit map
    \[
        \Emb(\Sigma,\RR^\infty)\longra \Emb(\Sigma,\RR^\infty)/\Diff(\Sigma)
    \]
    is identified as a principal $\Diff(\Sigma)$-bundle, and the quotient is denoted $B_\infty(\Sigma)$. Since $\Emb(\Sigma,\RR^\infty)$ is contractible by Whitney's embedding theorem in the stable Euclidean space, again as recalled in \cite[\S2.1, p.~4]{GMTW09}, this component is a model for $\mathbf{B}\!\Diff(\Sigma)$. Taking the coproduct over all diffeomorphism classes of closed $(d-1)$-manifolds proves condition~(\ref{ch3:: cond: bordism cat objects}).

    The proof of condition~(\ref{ch3:: cond: bordism cat morphisms}) is analogous, but the relevant family direction is now the simplicial variable $\AA^k$, not the interval direction. Fix $\sU(A,B)\in \cJ_1$, and for a $k$-simplex set
    \[
        W_{A,B}=W\cap \Big(\AA^k\times \big[\mathsf{mid}(A),\mathsf{mid}(B)\big]\times \RR^\infty\Big).
    \]
    Condition (\ref{id: spatial}) from \Cref{ch3:: def: bord_d} implies that the projection
    \[
        q_W \maps W_{A,B}\longra \AA^k
    \]
    is a submersion, and it is proper because $W$ is bounded and the interval $\big[\mathsf{mid}(A),\mathsf{mid}(B)\big]$ is compact. Therefore $W_{A,B}$ should be regarded as a smooth family, over $\AA^k$, of compact $d$-manifolds with boundary. Condition (\ref{id: categorical}), restricted to the distinguished closed pieces $A$ and $B$, identifies smoothly varying incoming and outgoing collars in this family. Hence each fiber of $q_W$ is naturally a compact bordism, and the whole family is a smoothly varying family of bordisms with fixed source and target.

    Since the base $\AA^k$ is contractible, this family is trivial up to isomorphism. After choosing one fiber $M$, the residual datum is again a smoothly varying family of neat embeddings of $M$ into $\big[\mathsf{mid}(A),\mathsf{mid}(B)\big]\times \RR^\infty$, well-defined only up to smooth maps $\AA^k\to \Diff(M,\{\partial_\nu M\})$. Thus each connected component is modeled on the quotient of a contractible neat embedding space by the collar-preserving diffeomorphism group, hence is a model for $\mathbf{B}\!\Diff(M,\{\partial_\nu M\})$. Taking the coproduct over diffeomorphism classes of bordisms proves condition~(\ref{ch3:: cond: bordism cat morphisms}).

    Lastly, let $\sU(I_0,\ldots,I_k)\in \cJ_k$. Here the same point applies: the relevant smooth family is again parametrized by the simplicial variable. A simplex of $\bord_d\big(\sU(I_0,\ldots,I_k)\big)$ determines, by restricting to the closed interval spanning the chosen midpoints, a proper submersion to $\AA^\bullet$ whose fibers are compact $d$-manifolds with boundary. Condition (\ref{id: categorical}) records, along the distinguished closed pieces of $\core\big(\sU(I_0,\ldots,I_k)\big)$, the collar neighborhoods of the intermediate cutting hypersurfaces. Repeating the same trivialization argument over the contractible base $\AA^\bullet$ shows that every connected component is determined by a fixed composed bordism
    \[
        \Sigma_0 \xra{M_1} \Sigma_1 \xra{M_2} \cdots \xra{M_k} \Sigma_k,
    \]
    together with the chosen collar neighborhoods of the intermediate hypersurfaces. The automorphism group of this data is precisely the group $\Diff(M,\{h_i\}_{i=0}^{k})$ from \Cref{ch3:: def: diffeo group}. As above, the remaining embedding data are contractible in the stable Euclidean ambient space, so the corresponding connected component is a model for $\mathbf{B}\!\Diff(M,\{h_i\}_{i=0}^{k})$. Taking the coproduct over all diffeomorphism classes of $k$-composable chains proves condition~(\ref{ch3:: cond: bordism cat composable}).

    We conclude that $\bord_d$ satisfies the three defining properties above, and hence it is the $(\oo,1)$-category of $d$-dimensional bordisms.
\end{proof}

\begin{remark}
    The proof above isolates the two geometric inputs that will need stratified analogues later on: Ehresmann's fibration theorem, and the triviality of smooth bundles over contractible bases. Together, these reduce each simplex to a smoothly varying family of embeddings of a fixed manifold, or bordism, into $\RR^\infty$.
\end{remark}








\section{Symmetric monoidal structure}

One of the defining features of the cobordism categories studied and used in topological field theories is that it is symmetric monoidal relative to disjoint union. Constructing a symmetric monoidal structure for an $(\oo,1)$-category can be done in different flavors, as discussed in \Cref{ch2:: sec: symmetric monoidal structure}. We opt for using $\Gamma$-spaces because it offers a clearer route to formalizing models of $(\oo,1)$-categories of symmetric monoidal $(\oo,1)$-categories.


\begin{construction}\label{ch3:: cons: Gamma-space}
  For every $m\geq 0$, let $\bord_d^{\langle m \rangle}$ denote the following interval Segal space:
\begin{align}
    \bord_d^{\langle m \rangle} \maps \Int &\longra \sset \\
    I &\longmapsto \Big(\,[k] \mapsto \{M \subset \AA^{k}\times \mathsf{cor}(I)\times \RR^{\oo}\}\,\Big)
\end{align} 
where $M$ is given by a disjoint union $M = M_{1} \sqcup \cdots \sqcup M_{m}$ where each $M_{i}\subset \AA^{k} \times \mathsf{cor}(I)\times \RR^{\oo}$ is a $k$-simplex of $\bord_d(I)$. Given a pointed map of pointed finite sets, $f\maps \langle m \rangle \longra \langle n \rangle$, we construct a map of interval Segal spaces. This is the natural transformation whose component on $I \in \Int$ is
\begin{equation}
    \begin{tikzcd}
        {\bord_d^{\langle m \rangle }(I)} & {M = M_1 \sqcup \cdots \sqcup M_m} \\
        \\
        {\bord_d^{\langle n \rangle }(I)} & {(f_*)_I = \Big(\coprod_{f(i) = 1}M_i\Big) \sqcup \cdots \sqcup \Big(\coprod_{f(i) = n}M_i\Big)}
        \arrow[from=1-1, to=3-1]
        \arrow["{{(f_*)_I}}"', maps to, from=1-2, to=3-2]
    \end{tikzcd}
\end{equation}
where we omit the ambient data of the embeddings for simplicity.
\end{construction}

\nid The structure outlined in \Cref{ch3:: cons: Gamma-space} can be assembled so as to endow $\bord_d$ with a structure of a symmetric monoidal interval Segal space, as defined in \Cref{ch2:: sec: symmetric monoidal structure}.

\begin{theorem}\label{ch3:: thm: symmetric monoidal}
    The assignment 
    \begin{align}
        \bord_{d}^{\langle \bullet \rangle} \maps \Gamma &\longra \mathsf{ISSp} \\
        \langle m \rangle & \longmapsto \bord_d^{\langle m \rangle}
    \end{align}
    is functorial and endows the interval Segal space $\bord_d$ of unoriented $d$-dimensional bordisms with a symmetric monoidal structure.
\end{theorem}

\begin{proof}

    Verification that $\bord_{d}^{\langle \bullet \rangle}$ is a symmetric monoidal structure reduces to proving that the induced $\Gamma$-Segal map
    \begin{equation}
        \prod_{1 \leq k \leq m} i_k \maps \bord_{d}^{\langle m \rangle} \longra \underbrace{\bord_{d}^{\langle 1 \rangle} \times \cdots \times \bord_{d}^{\langle 1 \rangle}}_{m \text{ times}}
    \end{equation}
    is an equivalence of interval Segal spaces, i.e., a level-wise equivalence — which we will henceforth denote as $S_m$, to avoid cluttering notation. To show this, fix embeddings
    \begin{equation}
        \iota_k \maps \RR^{\oo} \hookra \RR^{\oo}, \quad \iota_k(\RR^{\oo}) \cap \iota_l(\RR^{\oo}) = \emptyset \mbox{ for } k, l \in \{ 1, 2, \ldots, m \}, \, \, k \neq l .
    \end{equation}
    For every $I\in \Int$, we construct an explicit homotopy inverse
    \begin{equation}
        r_{m} \maps \bord_{d}^{\langle 1 \rangle} \times \cdots \times \bord_{d}^{\langle 1 \rangle} \xra{\,\,\,\,} \bord_{d}^{\langle m \rangle}
    \end{equation}
    where we first consider the disjoint union of all $M_i$ sitting inside $\AA^{k}\times \mathsf{cor}(I)\times \big(\RR^{\oo}\big)^{\sqcup m}$ and then take its image under the action of the coproduct of the $\iota_k$ for all $k=0,\ldots,m$; this is subsumed by the formula
    \begin{equation}
        (M_1,\ldots, M_m) \longmapsto \bigcup_{k = 1}^{m}(\id_{\AA^k\times (I \cup I)}\times \iota_k)(M_1\sqcup \cdots \sqcup M_m) \subset \AA^k \times \mathsf{cor}(I)\times \RR^{\oo}
    \end{equation}
    
    Both composites are shown to be homotopy equivalent to the identity using as argument the weak contractibility of the space $\Emb(\RR^{\oo},\RR^{\oo})$. Let us start with 
    \begin{equation}
         \Bigg(\prod_{0\leq k \leq m} i_k \Bigg)  \circ r_{m} \maps \bord_{d}^{\langle 1 \rangle}(I) \times \cdots \times \bord_{d}^{\langle 1 \rangle}(I) \longra \bord_{d}^{\langle 1 \rangle}(I) \times \cdots \times \bord_{d}^{\langle 1 \rangle}(I)  
    \end{equation}
    which sends a tuple of $k$-simplices $(M_1,\ldots,M_m)$ in $\bord_d(I)$ to the disjoint union post-composed with the embedding $M = M_1 \sqcup \cdots \sqcup M_m \subset \AA^{k}\times \mathsf{cor}(I)\times \RR^{\oo}$, before finally sending each $M_i$ to its corresponding $k$-simplex in $\bord_d^{\langle 1 \rangle}$ obtained by forgetting all the others. There is a homotopy equivalence from this map to the identity on $\bord_{d}^{\langle 1 \rangle} \times \cdots \times \bord_{d}^{\langle 1 \rangle}$ defined by constructing an isotopy $H\maps \RR^{\oo}\times \AA^1 \longra \RR^{\oo}$ which for the $k$-th entry sends $\iota_k$ to $\id{\RR^{\oo}}$. The isotopy is well-defined precisely because it acts solely on the $\RR^{\oo}$ factor, whilst preserving all submersivity conditions.

    The other composite
    \begin{equation}
          r_{m} \circ \Bigg(\prod_{0\leq k \leq m} i_k \Bigg) \maps \bord_{d}^{\langle m \rangle}(I) \longra \bord_{d}^{\langle m \rangle}(I)
    \end{equation}
    acts on a disjoint union $M = M_1 \sqcup \cdots \sqcup M_m$,  where $\big(M \subset \AA^{k} \times \mathsf{cor}(I)\times \RR^{\oo}\big)$ first by sending it to $(M_1,\ldots, M_m)$ where each $M_i$ is just $M$ with the other $(m-1)$ components deleted, followed by a re-embedding of each component into $\RR^{\oo}$ using the fixed choices of embeddings $\bigcup_{k = 0}^{m}\iota_k \maps (\RR^{\oo})^{\sqcup k} \hookra \RR^{\oo}$. This admits a homotopy from it to the identity on $\bord_{d}^{\langle m \rangle}(I)$ defined through an isotopy $H\maps \RR^{\oo}\times \AA^{1} \longra \RR^{\oo}$ chosen so that $H_1 \circ \iota_k = e_k$, where $e_k$ encodes the embedding of each $M_k$ into $\AA^{k}\times \mathsf{cor}(I)\times \RR^{\oo}$.   

    \begin{remark}[Operadic viewpoint]\label{ch43:: rem :operadic-symm-monoidal}
        It is worth noting that the symmetric monoidal structure on $\bord_d$ could also be packaged by exhibiting $\bord_d$ as an algebra over an $\cE_\infty$-operad, rather than via the Segal $\Gamma$-object formalism used here.

        Recall that one convenient model for $\cE_n$ is the \emph{little $n$-disks} (or little $n$-cubes) operad, whose $k$-ary operations are rectilinear embeddings
        \begin{equation}
        \bigsqcup_{i=1}^k (0,1)^n \hookrightarrow (0,1)^n
        \end{equation}
        (with operadic composition given by substitution of embeddings); see \cite[\S5.1.1]{LurieHA}.
        In particular, one may model $E_\infty$ by ``little cubes in $\RR^\infty$'', and algebras over such an operad are precisely commutative monoid objects up to coherent homotopy; cf.\ again \cite[\S5.1.1]{LurieHA}.

        In our geometric setting, for each rectilinear embedding
        \begin{equation}
        \iota:\;\bigsqcup_{i=1}^k \RR^\infty \hookrightarrow \RR^\infty
        \end{equation}
        disjoint union of $k$ families of bordisms (inside $M\times \core(T)\times \RR^\infty$) followed by postcomposition with $\iota$ defines a $k$-ary operation
        \begin{equation}
        \mu_\iota:\;\bord_{d}^{\times k}\longrightarrow \bord_{d}.
        \end{equation}
        The operadic compatibilities are inherited from functoriality of postcomposition with embeddings, and the contractibility of the relevant embedding spaces in $\RR^\infty$ forces the required higher coherences.

        We nevertheless work with the equivalent Segal $\Gamma$-object presentation, since it interfaces more directly with the model-categorical framework we have used so far.
    \end{remark}
\end{proof}





\section{Homotopy category as the bordism 1-category}

Before studying $(\oo,1)$-categories we studied ordinary 1-categories, and the first instances of the bordism category date back to at least Atiyah's 1988 seminal paper \cite{At88}. To confirm that we have successfully lifted the ordinary bordism category to the $(\oo,1)$-categorical setting, we should compute its associated symmetric monoidal  homotopy 1-category, and verify that they are equivalent as symmetric monoidal categories. To achieve this, we will use the construction presented in \Cref{ch2:: def: homotopy category ISSP}.


\begin{remark}\label{ch3:: rmk: homotopy ehresmann}
    In the next theorem we will work with the homotopy category of $\bord_{d}$. Recall from \Cref{ch2:: def: homotopy category ISSP} that this category is defined as the homotopy category of the Segal space 
    \[
    X  \,\overset{\sf def}{=}\, \mathsf{RKan}_{\gamma}(\bord_d).
    \]

    \nid The objects of this homotopy category are thus the 0-simplices of $X$. A 0-simplex of this homotopy limit can be understood as a coherent family $\big(\Sigma_{I}, [0] \xra{\id} \gamma(\sU(I))\big)$ indexed by objects of $\cJ_0$, where $\Sigma_{I}\in \bord_d\big(\sU(I))_0$. The coherence comes from the fact that for every $f\maps \sU(I) \ra \sU(J)$, there is a 1-simplex in $\bord_d\big(\sU(J)\big)$ from $(\Sigma_I)_{|\sU(J)}$; which amounts to a concordance $\widetilde{\Sigma}_{I,J}\in \bord_{d}(\sU(J))_1$ relating $\bord_{d}\big(\sU(J)\big)$ with $\Sigma_J$.
    
    \medskip

    \nid  We illustrate this with the following diagram, where we depict the pair 
    \medskip
    \begin{equation}
        \tikzfig{figures/Ch3/BOR15}{\begin{tikzpicture}
	\begin{pgfonlayer}{nodelayer}
		\node [style=none] (0) at (-2, 3) {};
		\node [style=none] (1) at (-2, 1) {};
		\node [style=none] (2) at (3, 3) {};
		\node [style=none] (3) at (3, 1) {};
		\node [style=none] (4) at (-5, -1) {};
		\node [style=none] (5) at (7, -1) {};
		\node [style=none] (6) at (-2, -1) {};
		\node [style=none] (7) at (3, -1) {};
		\node [style=none] (8) at (-10, -6) {};
		\node [style=none] (9) at (2, -6) {};
		\node [style=none] (10) at (-7, -6) {};
		\node [style=none] (11) at (-2, -6) {};
		\node [style=none] (12) at (-10, -7) {};
		\node [style=none] (13) at (-8, -7) {};
		\node [style=none] (14) at (-1, -7) {};
		\node [style=none] (15) at (2, -7) {};
		\node [style=none] (16) at (-7, -5) {};
		\node [style=none] (17) at (-2, -5) {};
		\node [style=none] (18) at (-7, -3) {};
		\node [style=none] (19) at (-7, -3) {};
		\node [style=none] (20) at (-7, -5) {};
		\node [style=none] (21) at (-2, -3) {};
		\node [style=none] (22) at (-2, -5) {};
		\node [style=none] (23) at (-2, -3) {};
		\node [style=none] (24) at (-7, -5) {};
		\node [style=none] (25) at (-8, -5) {};
		\node [style=none] (26) at (-8, -3) {};
		\node [style=none] (27) at (-1, -3) {};
		\node [style=none] (28) at (-1, -5) {};
		\node [style=none] (29) at (-5.25, -4.75) {};
		\node [style=none] (30) at (-3.75, -5.25) {};
		\node [style=none] (31) at (-5.75, -3.25) {};
		\node [style=none] (32) at (-4.5, -2.75) {};
		\node [style=none] (33) at (-3.75, -3.5) {};
		\node [style=none] (34) at (-11, -4) {$(\Sigma_{I})_{|\mathsf{U}(J)}$};
		\node [style=none] (35) at (-5, 2) {$\Sigma_{J}$};
		\node [style=none] (36) at (-5, -1) {};
		\node [style=none] (37) at (-2, -1) {};
		\node [style=none] (38) at (-10, -6) {};
		\node [style=none] (39) at (-7, -6) {};
		\node [style=endpoint] (40) at (-5, -1) {};
		\node [style=endpoint] (41) at (-10, -6) {};
		\node [style=endpoint] (42) at (2, -6) {};
		\node [style=endpoint] (43) at (7, -1) {};
		\node [style=none] (44) at (3, -1) {};
		\node [style=none] (45) at (7, -1) {};
		\node [style=none] (46) at (-2, -6) {};
		\node [style=none] (47) at (2, -6) {};
		\node [style=none] (48) at (3, 3) {};
		\node [style=none] (49) at (3, 1) {};
		\node [style=none] (50) at (-2, 1) {};
		\node [style=none] (51) at (-2, 3) {};
		\node [style=none] (52) at (-7, -5) {};
		\node [style=none] (53) at (-7, -3) {};
		\node [style=none] (54) at (-2, -3) {};
		\node [style=none] (55) at (-2, -5) {};
		\node [style=none] (56) at (-2, -3) {};
		\node [style=none] (57) at (-2, -5) {};
	\end{pgfonlayer}
	\begin{pgfonlayer}{edgelayer}
		\draw [style=geometric] (12.center) to (15.center);
		\draw [style=geometric] (0.center) to (2.center);
		\draw [style=geometric] (1.center) to (3.center);
		\draw [style=geometric] (4.center) to (5.center);
		\draw [style={open_int}] (4.center) to (6.center);
		\draw [style={open_int}] (7.center) to (5.center);
		\draw [style=geometric] (8.center) to (9.center);
		\draw [style={open_int}] (8.center) to (10.center);
		\draw [style={open_int}] (11.center) to (9.center);
		\draw [style={open_int_blue}] (12.center) to (13.center);
		\draw [style={open_int_blue}] (14.center) to (15.center);
		\draw [style=ts-dashed-grey] (25.center) to (24.center);
		\draw [style=ts-dashed-grey] (26.center) to (19.center);
		\draw [style=ts-dashed-grey, in=-180, out=-180, looseness=0.75] (26.center) to (25.center);
		\draw [style=ts-dashed-grey, in=0, out=0, looseness=0.75] (26.center) to (25.center);
		\draw [style=ts-dashed-grey] (22.center) to (28.center);
		\draw [style=ts-dashed-grey] (23.center) to (27.center);
		\draw [style=ts-dashed-grey, in=180, out=-180, looseness=0.75] (27.center) to (28.center);
		\draw [style=ts-dashed-grey, in=0, out=0, looseness=0.75] (27.center) to (28.center);
		\draw [style=geometric, in=210, out=0] (24.center) to (29.center);
		\draw [style=geometric, in=-180, out=0, looseness=1.25] (30.center) to (22.center);
		\draw [style=geometric, in=180, out=0] (29.center) to (30.center);
		\draw [style=geometric, in=165, out=0] (19.center) to (31.center);
		\draw [style=geometric, in=180, out=0] (31.center) to (32.center);
		\draw [style=geometric, in=180, out=0] (32.center) to (33.center);
		\draw [style=geometric, in=180, out=0] (33.center) to (23.center);
		\draw [style={open_int_blue_2_dash}] (25.center) to (13.center);
		\draw [style={open_int_blue_2_dash}] (28.center) to (14.center);
		\draw [style=defect line 3 dashed] (24.center) to (10.center);
		\draw [style=defect line 3 dashed] (22.center) to (11.center);
		\draw [style=defect line 3 dashed] (1.center) to (6.center);
		\draw [style=ts-dashed-red] (3.center) to (7.center);
		\draw [style=red shade low op] (36.center)
			 to (38.center)
			 to (39.center)
			 to (37.center)
			 to cycle;
		\draw [style=red shade low op] (47.center)
			 to (46.center)
			 to (44.center)
			 to (45.center)
			 to cycle;
		\draw [style=geometric, in=-180, out=180, looseness=0.75] (48.center) to (49.center);
		\draw [style=geometric, in=0, out=0, looseness=0.75] (48.center) to (49.center);
		\draw [style=geometric, in=180, out=-180, looseness=0.75] (50.center) to (51.center);
		\draw [style=ts-dashed, in=0, out=0, looseness=0.75] (51.center) to (50.center);
		\draw [style=geometric, in=180, out=-180, looseness=0.75] (52.center) to (53.center);
		\draw [style=ts-dashed, in=0, out=0, looseness=0.75] (53.center) to (52.center);
		\draw [style=geometric, in=-180, out=180, looseness=0.75] (56.center) to (57.center);
		\draw [style=geometric, in=0, out=0, looseness=0.75] (56.center) to (57.center);
		\draw [style=ts-dashed-grey] (51.center) to (53.center);
		\draw [style=ts-dashed-grey] (52.center) to (50.center);
		\draw [style=ts-dashed-grey] (48.center) to (56.center);
		\draw [style=ts-dashed-grey] (57.center) to (49.center);
	\end{pgfonlayer}
\end{tikzpicture}
}
    \end{equation} 

    \medskip
    
    \nid The bottom line is that, for every $0$-simplex and its associated coherent family, we can associate a unique diffeomorphism class of $(d-1)$-dimensional unoriented smooth manifolds. This class is given by considering the diffeomorphism class of the fiber of 
    \begin{equation}
    \Sigma_I \longra \sU(I) \times \RR^{\oo}
    \end{equation}
    over any $t\in \sU(I)$. This fiber is constant, up to diffeomorphism, since it is a smooth fiber bundle as a consequence of \Cref{ap:: prop: ehresmann boundary}.
\end{remark}


\begin{theorem}\label{ch3:: thm: homotopy category bordism}
    There is an equivalence of symmetric monoidal 1-categories 
    \begin{equation}
        \ho(\bord_d) \,\, \simeq \,\, \mathsf{Cob}_d
    \end{equation}
    between the homotopy category of $\bord_{d}$ and the unoriented $d$-dimensional bordism category.
\end{theorem}

\begin{proof}
    We construct functor 
    \begin{equation}
        h \maps \ho(\bord_d) \longra \mathsf{Cob}_d.
    \end{equation}
    Let $\underline{\Sigma} = (\Sigma_{I}, [0] \xra{\id} \gamma(I),\pi_{I})$ denote an object of $\ho(\bord_d)$, i.e., a coherent family of $(d-1)$-dimensional smooth unoriented manifolds. We define
    \begin{equation}
        h \big(\underline{\Sigma}\big) \, \overset{\mathsf{def}}{=} \, \pi_{I}^{-1}(\mathsf{mid}(I)), \quad \text{for some } I\in \cJ_0.
    \end{equation}
    By \Cref{ch3:: rmk: homotopy ehresmann}, this is well-defined up to diffeomorphism.

    Now fix objects $\underline{\Sigma}$ and $\underline{\Sigma'}$ of $\ho(\bord_d)$. A morphism in
    \[
        \Hom_{\ho(\bord_d)}\big(\underline{\Sigma},\underline{\Sigma'}\big)
    \]
    is represented by a coherent family of $d$-dimensional unoriented bordisms compatible with the given source and target data. Choosing some $\sU(A,B)\in \cJ_1$ and a representative
    \[
        M \hookra \sU(A,B)\times \RR^{\oo},
    \]
    we define
    \begin{equation}
        M_{A,B} = \pi^{-1}\big([\mathsf{mid}(A),\mathsf{mid}(B)]\big), \quad h\big([M]\big) \, \overset{\sf def}{=} \, \big[M_{A,B}\big],
    \end{equation}
    where we take the diffeomorphism class of the extracted bordism $M_{A,B}$. This is well-defined: changing the choice of $\sU(A,B)$ only changes the extracted fiber by a diffeomorphism, again by \Cref{ch3:: rmk: homotopy ehresmann}, and concordant representatives determine the same class in $\mathsf{Cob}_d$.

    \medskip

    To prove that this functor is essentially surjective, let $\Sigma$ be a closed $(d-1)$-dimensional smooth manifold. The cylinder $\Sigma\times [0,1]$ admits a neat embedding into $\sU(A)\times \RR^{\oo}$ for every $\sU(A)\in \cJ_{0}$ by the manifold-with-boundary version of Whitney's embedding theorem, for instance \cite[Theorem 8.9]{CS19}. Choosing these embeddings compatibly, via linear diffeomorphisms $\sU(A)\cong [0,1]$, they assemble into a coherent family and hence define an object of $\ho(\bord_d)$ whose image under $h$ is the diffeomorphism class of $\Sigma$.

    \medskip
    
    It remains to show that $h$ is fully faithful. Fix objects $\underline{\Sigma}$ and $\underline{\Sigma'}$ of $\ho(\bord_d)$, and write
    \[
        \Sigma = h(\underline{\Sigma}), \qquad \Sigma' = h(\underline{\Sigma'}).
    \]
    By \Cref{ch3:: thm: bordism category characterization}, the total space of $1$-morphisms in $\bord_d$ is weakly equivalent to a coproduct of classifying spaces indexed by diffeomorphism classes of bordisms. Restricting the source-target map to the fiber over the pair $(\underline{\Sigma},\underline{\Sigma'})$ therefore selects precisely the sub-coproduct indexed by bordisms from $\Sigma$ to $\Sigma'$. Hence the corresponding mapping space is weakly equivalent to
    \begin{equation}
        \coprod_{[M \maps \Sigma \longra \Sigma']} \mathbf{B}\!\Diff(M,\{\partial_{\nu}M\}),
    \end{equation}
    where the coproduct ranges over diffeomorphism classes of compact $d$-dimensional bordisms from $\Sigma$ to $\Sigma'$. Passing to path components, and using that each classifying space $\mathbf{B}\!\Diff(M,\{\partial_{\nu}M\})$ is path connected, we obtain a bijection
    \begin{equation}
        \Hom_{\ho(\bord_d)}\big(\underline{\Sigma},\underline{\Sigma'}\big)
        \cong
        \big\{[M \maps \Sigma \longra \Sigma']\big\}
        =
        \Hom_{\mathsf{Cob}_d}\big(h(\underline{\Sigma}),h(\underline{\Sigma'})\big).
    \end{equation}
    Under this identification, the induced map on hom-sets is exactly the map defined above by extracting the bordism over the interval between the chosen midpoints. Hence $h$ is fully faithful. Since it is also essentially surjective, it is an equivalence of categories. Compatibility with disjoint union is immediate from the construction, so this is an equivalence of symmetric monoidal $1$-categories.
\end{proof}

\section{Tangential structures}
\label{ch3:: sec: tangential structures}

When studying manifold topology, it is natural to ask what additional structures manifolds may be endowed with. Generally speaking, the structures we care about are the simplest ones: \emph{tangential structures}. These structures can be thought of as families of linear data, such as group representations.

\begin{definition}
    An $n$-dimensional tangential structure is a topological space $B$ together with a fibration $B \longra \BO(n)$. A \define{stable tangential structure} is a topological space $\fB$ together with a fibration $\fB \longra \BO$. 
\end{definition}

\begin{remark}
    A stable tangential structure $\fB \longra \BO$ induces, for every $n\geq 0$, an $n$-dimensional tangential structure via pullback 
    \begin{equation}
        \begin{tikzcd}[ampersand replacement=\&]
            {B(n)} \&\& {\BO(n)} \\
            \\
            \fB \&\& \BO
            \arrow[from=1-1, to=1-3]
            \arrow[from=1-1, to=3-1]
            \arrow["\lrcorner"{anchor=center, pos=0.125}, draw=none, from=1-1, to=3-3]
            \arrow[hook, from=1-3, to=3-3]
            \arrow[from=3-1, to=3-3]
        \end{tikzcd}
    \end{equation}
    where $\BO(n) \hookra \BO$ denotes the canonical inclusion.
\end{remark}

\begin{remark}
    Since we can always factor a map of topological spaces $B \longra \BO(n)$ as \[B \xra{\,\,\,\sim\,\,\,} B^{\sf fib} \longra \BO(n),\] where the second map is a fibration, we might as well define an $n$-dimensional tangential structure as a map $B \longra \BO(n)$.
\end{remark}

\begin{definition}
    Let $M$ be a $k$-dimensional smooth manifold and let $B\longra \BO(n)$ be an $n$-dimensional tangential structure. A $B$\define{-framing} on $M$ is a map $\alpha \maps M \longra B$, together with a homotopy making the following diagram commute up to the specified homotopy:
    \begin{equation}
    \begin{tikzcd}[ampersand replacement=\&]
        \&\& B \\
        M \&\& {\BO(n).}
        \arrow[from=1-3, to=2-3]
        \arrow["\alpha", from=2-1, to=1-3]
        \arrow[""{name=0, anchor=center, inner sep=0}, "{\widetilde{\mathsf{\tau}_M}}"', from=2-1, to=2-3]
        \arrow[between={0.4}{0.6}, Rightarrow, from=1-3, to=0]
    \end{tikzcd}
    \end{equation}
    The map $\widetilde{\tau}_M$ is the map classifying the $n$-dimensional stabilized tangent bundle 
    \begin{equation}
        \widetilde{\sT M} \,\, \overset{\sf def}{=} \,\, \epsilon^{n-k}_{M} \oplus \sT M.
    \end{equation}   
\end{definition}

\begin{remark}
    Instead of defining a $B$-framing as a lift together with a homotopy, we can equivalently state that it is a lift in $\Top$ as an $(\oo,1)$-category — this makes the choice of a homotopy implicit.
\end{remark}

\begin{example}
    Let $G$ be a Lie group together with a linear representation $\rho \maps G \ra \sG \sL(n$). Recall that the Gram-Schmidt orthogonalization offers a strong deformation retract of $\sO(n) \xhookrightarrow{\simeq} \sG\sL(n)$. The representation $\rho$ gives rise to a map of spaces
    \begin{equation}
        \rho \maps \sB G \longra \sB G\sL(n) \xra{\simeq} \sB\sO(n).
    \end{equation}
    This representation classifies a smooth rank $n$ bundle $E$ on $\sB G$. To endow a smooth $n$-dimensional manifold $M$ with a $(\sB G,\rho)$-framing is to give a map $\varphi\maps M \longra \sB\sG$ together with an isomorphism of rank $n$ vector bundles over $M$:
    \begin{equation}
        \varphi^{*}E \xra{\,\,\,\,\cong\,\,\,\,} \sT M.
    \end{equation}
    Ultimately, a $\sB G$-framing amounts to a reduction of the structure group of $\sT M$ from $\sO(n)$ to $\SO(n)$. Equivalently, a manifold admits a $\sB G$-framing if and only if its tangent bundle's transition functions can be chosen so as to take values in $G$.
\end{example}

\begin{example}
    The short exact sequence 
    \begin{equation}
        1 \longra \SO (n) \longra \sO(n) \longra \sO(1) \longra 1
    \end{equation}
    induces a group homomorphism $\SO(n) \longra \sO (n)$ for every $n \geq 1$. A smooth $n$-manifold with a $\sB\!\SO (n)$-framing is thus a smooth manifolds endowed with an additional choice of an orientation.   
\end{example}

\begin{example}
    Not every tangential structure is obtained from a smooth linear representation of a Lie group. Let $X$ be a space, an $X$-framing is defined via the projection 
    \begin{equation}
        X \times \sB\sO(n) \xra{\sf pr} \sB\sO(n).
    \end{equation}
    An $X$-framed smooth manifold $M$ is a smooth $n$-dimensional manifold together with a continuous map $M \longra X \times \sB\sO(n)$ over $\sB\sO(n)$; this data is determined by a continuous map $M \longra X$. In the case where $X = \cD_n$ is a discrete space, i.e., a set, a continuous map $M \longra \cD_n$ factors through $M \ra \pi_{0}(M)$, which is equivalent to a \emph{coloring} or a labeling of the connected components of $M$ by colors, the elements of $\cD_n$. This tangential structure admits immediate generalizations: we can combine a smooth linear representation of a Lie group $G$, together with a coloring by $\cD_n$:
    \begin{equation}
        \big(B \longra \sB\sO(n)\big) \simeq \big(\cD_{n}\times \sB G \xra{\,\,\, \id \times \rho\,\,\,} \cD_{n}\times \sB \sO(n) \xra{\,\,\,\pr\,\,\,} \sB\sO(n)\big).
    \end{equation}
    An $(\cD_n\times G,\rho)$-framing on $M$ amounts to a coloring of its connected components by colors in $\cD_n$ together with a $G$-reduction of its tangent bundle.
\end{example}

When we have a family of smooth manifolds, i.e. a smooth bundle of manifolds with fiber a smooth manifold $M$, it is natural to ask how might families of smooth manifolds with a given tangential structure behave. To study such families, we introduce the \emph{fiberwise tangent bundle of a smooth bundle $\pi_W\maps W \longra M$} as the kernel of the derivative of $\pi_W$:

\begin{equation}
    \mathsf{Ker}(\sD\pi_{l}) \longra M.
\end{equation}
This bundle is also often called the vertical tangent bundle.

Using the fiberwise tangent bundle we can formulate what we mean by a family of $B$-framed smooth manifolds. Roughly, a family of $B$-framed smooth manifolds is a smooth fiber bundle $\pi_W \maps W \longra M$ together with a $B$-trivialization of its fiberwise tangent bundle.

\begin{definition}
    Let $\pi_W \maps W \longra M$ be a smooth fiber bundle. We say that it is a smooth family of $B$-framed manifolds if there is a map $\alpha\maps M \longra B$ together with a homotopy making the following diagram commute:
    \begin{equation}
    \begin{tikzcd}[ampersand replacement=\&]
        \&\& B \\
        M \&\& {\BO(n).}
        \arrow[from=1-3, to=2-3]
        \arrow["\alpha", from=2-1, to=1-3]
        \arrow[""{name=0, anchor=center, inner sep=0}, "{\widetilde{\mathsf{\tau}_{W/M}}}"', from=2-1, to=2-3]
        \arrow[between={0.4}{0.6}, Rightarrow, from=1-3, to=0]
    \end{tikzcd}
    \end{equation}
    The map $\widetilde{\mathsf{\tau}_{W/M}}\maps M \longra \BO(n)$ is the classifying map of the fiberwise tangent bundle over $M$.
\end{definition}

\begin{definition}\label{ch3:: def: bord_d tangential}
    Let $V$ be a finite-dimensional real vector space and let $e\maps B \longra \sB\sO(d)$ be a map of spaces. Let us define a functor 
    \begin{equation}
    \widetilde{\bord}^{B}_{d,V} \maps \mathsf{Man^{\op}}\times \Int \longrightarrow \set    
    \end{equation}
    sending a pair $(M,T)$, where $M\in \man^{\op}$ and $T\in \Int$, to the set of bounded closed subsets 
    \begin{equation}
    W\subseteq M\times \mathsf{cor}(T) \times V
    \end{equation} 
    satisfying the following four conditions:
    \begin{enumerate}
    \item $W$ is a smooth $(d+\mathsf{dim}(M))$-dimensional submanifold with boundary $\partial W$ such that 
    \begin{equation}
        \partial W \subset \partial \big(M \times \mathsf{cor}(T) \times V\big) = M \times \partial (\mathsf{cor}(T)) \times V   
    \end{equation}
    the boundary of $W$ must be contained in the boundary of the ambient space.
    \item The restricted projection 
    \begin{equation}
        \pi\maps \big(W \cap (M\times T^{\mathsf c} \times V) \big) \hookra  M \times T^{\mathsf c} \times V \fib  M \times T^{\mathsf c}
    \end{equation} 
    is a submersion. 
    \item The restricted projection 
    \begin{equation}
    \pi_W\maps W \hookra M\times \mathsf{cor}(T) \times V \fib M 
    \end{equation}
    is a submersion.
    \item A map of spaces $\alpha \maps W \longra B$ making the following diagram commute up to homotopy:
    \begin{equation}
        \begin{tikzcd}[ampersand replacement=\&]
            \&\& B \\
            W \&\& {\BO(d)}
            \arrow["e", from=1-3, to=2-3]
            \arrow["\alpha", dashed, from=2-1, to=1-3]
            \arrow["{\widetilde{\tau_{W/M}}}"', from=2-1, to=2-3]
        \end{tikzcd}
    \end{equation}   
    This amounts to a choice of a map together with a homotopy from $e\circ \alpha$ to $\widetilde{\tau}_{W/M}$, where the latter classifies the fiberwise tangent bundle of the smooth submersion $\pi_W\maps W \longra M$. 
    \end{enumerate} 

    \noindent The functor $\widetilde{\bord}^{B}_{d,V}$ sends:
    \begin{itemize}
    \item[($\man$):] The morphism $(f,\mathsf{id}_{T})\maps (N,T) \longra (M,T)$ induced by the smooth map $f\maps N \longra M$, which sends the submanifold with boundary
    \begin{equation}
         W\subseteq M\times \mathsf{cor}(T) \times V 
    \end{equation}
    to $(f \times \id)^{*} W$ — the pullback submanifold 
    \begin{equation}
        \begin{tikzcd}[ampersand replacement=\&]
            {(f\times \id)^\ast W} \& {W} \\
            \\
            N\times \mathsf{cor}(T) \times V \& M\times \mathsf{cor}(T) \times V,
            \arrow[from=1-1, to=1-2]
            \arrow[hook, "(f\times \id)^\ast \pi_{|W}"',from=1-1, to=3-1]
            \arrow["\lrcorner"{anchor=center, pos=0.125}, draw=none, from=1-1, to=3-2]
            \arrow[hook, from=1-2, to=3-2]
            \arrow["f\times \id"', from=3-1, to=3-2]
        \end{tikzcd}
    \end{equation}
    which is a $\big(d + \mathsf{dim}(N)\big)$-dimensional submanifold with boundary
    \begin{equation}
        (f\times \id)^* W \subseteq N \times \mathsf{cor}(T) \times V.
    \end{equation}
Which is well-defined since the pullback preserves the fibers over the boundary $\partial (N \times \mathsf{cor}(T) \times V)$. Moreover, we have that $(f\times \id)^* W$ comes equipped with a map of spaces to $\sB\sO(d)$ and a corresponding homotopy, both obtained through pre-composition with the top map of the pullback diagram
    \begin{equation}
        \begin{tikzcd}
            &&& B \\
            \\
            {(f\times \id)^*W} & W && {\sB\sO(n)}
            \arrow[from=1-4, to=3-4]
            \arrow[from=3-1, to=3-2]
            \arrow[dashed, from=3-2, to=1-4]
            \arrow["{\sT_{W/M}}"', from=3-2, to=3-4]
        \end{tikzcd}
    \end{equation}

    \item[($\Int$):] The pair $(\id_M, T\subset S)$ by \emph{forgetting} connected components and restricting the submanifold with boundary 
    \begin{equation}
        W \subseteq N \times \mathsf{cor}(S) \times V
    \end{equation}
    along $N \times \mathsf{cor}(T)\times V \hookra N \times \mathsf{cor}(S)\times V$. The restriction $W_{|T}$ comes equipped with a smooth embedding 
    \begin{equation}
        W_{|T} \hookra W
    \end{equation}
    which induces a $B$-framing on the restricted fiberwise tangent bundle, as desired.
    \end{itemize}  
\end{definition}


\begin{remark}
    The proof of the interval Segal condition rests entirely on the two locality statements \Cref{ch3:: prop: sheaf condition} and \Cref{ch3:: prop: naive sheaf con}. These also hold for $\widetilde{\bord}_{d}^B$ since $B$-framed manifolds enjoy the same locality properties. Therefore, we skip the proof of the next proposition.
\end{remark}

\begin{proposition}\label{ch3:: prop: Segal condition with structrue}
    Let $d \geq 0$ be a non-negative integer, let $B \longra \BO(d)$ be a tangential structure, and let $V$ be a finite-dimensional vector space of sufficiently high dimension. The composite functor
    \begin{align}
        \bord_{d,V}^{B} \maps \Int &\longra \shv(\man)  \longra \sset  \\
         I &\longmapsto \widetilde{\bord}_{d,V}^{B}(-,I) \longmapsto \bB\widetilde{\bord}_{d,V}^{B}(*,I) = \left|\,\widetilde{\bord}_{d,V}^{B}(\AA^\bullet,I)\,\right| 
    \end{align}
    is an interval Segal space. 
\end{proposition}










